% ============ Chapter 5. Dimension theory (原书 pp.92-123) ============
\chapter{Dimension theory}\label{ch:5}

The dimension theory of Noetherian rings is probably the greatest of
Krull's many achievements; with his principal ideal theorem
(Theorem~\ref{thm:13.5}) the theory of Noetherian rings gained in mathematical
profundity. Then the theory of multiplicities was first treated rigorously
and in considerable generality by Chevalley, and was simplified by Samuel's
definition of multiplicity in terms of the Samuel function.

Here we follow the method of EGA, proving Theorem~\ref{thm:13.4} via the
Samuel function, and deducing the principal ideal theorem as a corollary.
The Samuel function is of importance as a measure of singularity in
Hironaka's resolution of singularities, but in this book we can only cover its
basic properties. In \S\ref{sec:15} we exploit the notion of systems of parameters to
discuss among other things the dimension of the fibres of a ring
homomorphism and the dimension formula for finitely generated extension
rings.

% ============ §13 Graded rings, the Hilbert function and the Samuel function (原书 pp.92-101) ============
\section{Graded rings, the Hilbert function and the Samuel function}\label{sec:13}\index{graded rings, section 13}\index{Hilbert function, section 13}\index{Samuel function, section 13}

Let $G$ be an Abelian semigroup with identity element $0$; (that is, $G$
is a set with an addition law $+$ satisfying associativity $(x + y) + z = x +
(y + z)$, commutativity $x + y = y + x$, and such that $0 + x = x$). A \emph{graded} (or
$G$-graded) ring is a ring $R$ together with a direct sum decomposition of $R$ as
an additive group $R = \bigoplus_{i \in G} R_i$ satisfying $R_iR_j \subset R_{i+j}$. Similarly, a \emph{graded}
$R$-module is an $R$-module $M$ together with a direct sum decomposition $M =
\bigoplus_{i \in G} M_i$ satisfying $R_iM_j \subset M_{i+j}$. An element $x \in M$ is \emph{homogeneous}\index{homogeneous element} if
$x \in M_i$ for some $i \in G$, and $i$ is then called the \emph{degree} of $x$. A general element
$x \in M$ can be written uniquely in the form $x = \sum_{i \in G} x_i$ with $x_i \in M_i$ and only
finitely many $x_i \neq 0$; $x_i$ is called the \emph{homogeneous term} of $x$ of degree $i$.

A submodule $N \subset M$ is called a \emph{homogeneous submodule}\index{homogeneous submodule} (or \emph{graded}
submodule) if it can be generated by homogeneous elements. This
condition is equivalent to either of the following two:
\begin{enumerate}[(1)]
  \item For $x \in M$, if $x \in N$ then each homogeneous term of $x$ is in $N$;
\end{enumerate}
For a homogeneous submodule $N \subset M$ we set $N_i = N \cap M_i$; then
$M/N = \bigoplus_{i \in G} M_i/N_i$ is again a graded $R$-module.

One sees from the definition that $R_0 \subset R$ is a subring, and that each
graded piece $M_i$ of a graded $R$-module $M$ is an $R_0$-module.

The notion of graded ring is most frequently used when $G$ is the
semigroup $\{0, 1, 2, \ldots\}$ of non-negative integers, which we denote by $\N$. In this
case, we set $R^+ = \bigoplus_{n > 0} R_n$; then $R^+$ is an ideal of $R$, with $R/R^+ \cong R_0$.

The polynomial ring $R = R_0[X_1, \ldots, X_r]$ over a ring $R_0$ is usually made
into an $\N$-graded ring by defining the degree of a monomial $X_1^{\alpha_1} \cdots X_r^{\alpha_r}$ as
the total degree $\alpha_1 + \cdots + \alpha_r$; however, $R$ has other useful gradings. For
example, $R$ has an $\N^r$-grading in which $X_1^{\alpha_1} \cdots X_r^{\alpha_r}$ has degree $(\alpha_1, \ldots, \alpha_r)$;
the value of systematically using this grading can be seen in Goto--
Watanabe [1]. Alternatively, giving each of the $X_i$ some suitable weight
$d_i$ and letting the monomial $X_1^{\alpha_1} \cdots X_r^{\alpha_r}$ have weight $\sum \alpha_i d_i$ defines an
$\N$-grading of $R$. For example, the ring $R_0[X, Y, Z]/(f)$, where $f = a_1X^{\alpha_1} +
a_2Y^{\alpha_2} + a_3Z^{\alpha_3}$ can be graded by giving the images of $X, Y, Z$ the weights
$\alpha_1, \alpha_2$ and $\alpha_3$, respectively.

A \emph{filtration}\index{filtration} of a ring $A$ is a descending chain $A = J_0 \supset J_1 \supset \cdots$ of ideals
such that $J_nJ_m \subset J_{n+m}$; the associated graded ring $\gr(A)$ is defined as
follows. First of all as a module we set $\gr_n(A) = J_n/J_{n+1}$ for $n \ge 0$, and
$\gr(A) = \bigoplus_{n \ge 0} \gr_n(A)$; then we define the product by
\[
  (x + J_{n+1})(y + J_{m+1}) = xy + J_{n+m+1} \quad \text{for } x \in J_n \text{ and } y \in J_m.
\]
It is easy to see that $\gr(A)$ becomes a graded ring. The filtration
$J_1 \supset J_2 \supset \cdots$ defines a linear topology on $A$ (see \S\ref{sec:8}), and the completion $\hat{A}$ of
$A$ in this topology has a filtration $\hat{J}_1 \supset \hat{J}_2 \supset \cdots$ such that $\hat{A}/\hat{J}_n \cong A/J_n$ for
all $n$, hence $\hat{J}_n/\hat{J}_{n+1} \cong J_n/J_{n+1}$, and
\[
  \gr(\hat{A}) = \gr(A).
\]

Let $A$ be a ring, $I$ an ideal, and let $B = \bigoplus_{n \ge 0} I^n/I^{n+1}$ be the graded ring
associated with the filtration $I \supset I^2 \supset \cdots$ of $A$ by powers of $I$; the various
notations $\gr_I(A)$, $\gr^I(A)$ and $G_I(I)$ are used to denote $B$ in the current
literature. An element of $B_n = I^n/I^{n+1}$ can be expressed as a linear
combination of products of $n$ elements of $B_1 = I/I^2$, so that $B$ is generated
over the subring $B_0 = A/I$ by elements of $B_1$. If $I = Ax_1 + \cdots + Ax_r$ and $\bar{x}_i$
denotes the image of $x_i$ in $B_1 = I/I^2$ then
\[
  B = \gr_I(A) = (A/I)[\bar{x}_1, \ldots, \bar{x}_r],
\]
and $B$ is a quotient of the polynomial ring $(A/I)[X_1, \ldots, X_r]$ as a graded
ring.

\begin{theorem}\label{thm:13.1}
An $\N$-graded ring $R = \bigoplus_{n \ge 0} R_n$ is Noetherian if and only if
$R_0$ is Noetherian and $R$ is finitely generated as a ring over $R_0$.
\end{theorem}

\begin{proof}
The `if' is obvious, and we prove the `only if': suppose that $R$
is Noetherian. Then since $R_0 \cong R/R^+$, $R_0$ is Noetherian. $R^+$ is a home-
geneous ideal, and is finitely generated, so that we can suppose that it is
generated by homogeneous elements $x_1, \ldots, x_r$. Then it is easy to see that
$R = R_0[x_1, \ldots, x_r]$; in fact it is enough to show that $R_n \subset R_0[x_1, \ldots, x_r]$
for every $n$. Now writing $d_i$ for the degree of $x_i$ we have
\[
  (*) \quad R_n = x_1R_{n-d_1} + x_2R_{n-d_2} + \cdots + x_rR_{n-d_r}.
\]
Indeed, for $y \in R_n$, write $y = \sum c_ix_i$ with $f_i \in R$; then setting $g_i$ for the
homogeneous term of degree $n - d_i$ of $f_i$ (with $g_i = 0$ if $n - d_i < 0$), we also
have $y = \sum x_ig_i$. From $(*)$ it follows by induction that $R_n \subset R_0[x_1, \ldots, x_r]$.
\end{proof}

Let $R = \bigoplus_{n \ge 0} R_n$ be a Noetherian graded ring; then if $M = \bigoplus_{n \ge 0} M_n$
is a finitely generated graded $R$-module, each $M_n$ is finitely generated
as $R_0$-module. In fact when $M = R$ this is clear from $(*)$ above. In the general
case $M$ can be generated by a finite number of homogeneous elements $\omega_i$: $M
= \sum R\omega_i$. Now letting $e_i$ be the degree of $\omega_i$, we have as above that
\[
  M_n = R_{n-e_1}\omega_1 + \cdots + R_{n-e_r}\omega_r \quad \text{(where } R_i = 0 \text{ for } i < 0),
\]
and hence $M_n$ is a finite $R_0$-module. In particular if $R_0$ is an Artinian ring,
then $l(M_n) < \infty$, where $l$ denotes the length of an $R_0$-module. In this case we
define the \emph{Hilbert series}\index{Hilbert series} $P(M, t)$ of $M$ by the formula:
\[
  P(M, t) = \sum_{n=0}^{\infty} l(M_n)t^n \in \Z[[t]].
\]
(In combinatorics it is a standard procedure to associate with a sequence of
numbers $a_0, a_1, a_2, \ldots$ the generating function $\sum a_it^i$.)

\begin{theorem}\label{thm:13.2}
Let $R = \bigoplus_{n \ge 0} R_n$ be a Noetherian graded ring with $R_0$
Artinian, and let $M$ be a finitely generated graded $R$-module. Suppose that
$R = R_0[x_1, \ldots, x_r]$ with $x_i$ of degree $d_i$, and that $P(M, t)$ is as above. Then
$P(M, t)$ is a rational function of $t$, and can be written
\[
  P(M, t) = f(t) \Big/ \prod_{i=1}^{r} (1 - t^{d_i}),
\]
where $f(t)$ is a polynomial with coefficients in $\Z$.
\end{theorem}

\begin{proof}
By induction on $r$, the number of generators of $R$. When $r = 0$ we
have $R = R_0$, so that for $n$ sufficiently large, $M_n = 0$, and the power series
$P(M, t)$ is a polynomial. When $r > 0$, multiplication by $x_r$ defines an $R_0$-
linear map $M_n \to M_{n+d_r}$; writing $K_n$ and $L_{n+d_r}$ for the kernel and cokernel,
we get an exact sequence
\[
  0 \to K_n \to M_n \to M_{n+d_r} \to L_{n+d_r} \to 0.
\]
Set $K = \bigoplus K_n$, and $L = \bigoplus L_n$. Then $K$ is a submodule of $M$, and $L = M/x_rM$,
so that $K$ and $L$ are finite $R$-modules; moreover $x_rK = x_rL = 0$, so that $K$
and $L$ can be viewed as $R/x_rR$-modules, and hence we can apply the
induction hypothesis to $P(K, t)$ and $P(L, t)$. Now from the above exact
sequence we get
\[
  l(K_n) - l(M_n) + l(M_{n+d_r}) - l(L_{n+d_r}) = 0.
\]
If we multiply this by $t^{n+d_r}$ and sum over $n$ this gives
\[
  t^{d_r}P(K, t) - t^{d_r}P(M, t) + P(M, t) - P(L, t) = g(t),
\]
where $g(t) \in \Z[t]$. The theorem follows at once from this.
\end{proof}

A lot of information on the values of $l(M_n)$ can be obtained from the
above theorem. Especially simple is the case $d_1 = \cdots = d_r = 1$, so that $R$ is
generated over $R_0$ by elements of degree $1$. In this case $P(M, t) = f(t)
(1 - t)^{-r}$; if $f(t)$ has $(1 - t)$ as a factor we can cancel to get $P$ in the form
\[
  P(M, t) = f(t)(1 - t)^{-d} \quad \text{with } f \in \Z[t], \; d \ge 0,
\]
and if $d > 0$ then $f(1) \neq 0$.
If this holds, we will write $d = d(M)$. Since $(1 - t)^{-1} = 1 + t + t^2 + \cdots$, we can
repeatedly differentiate both sides to get
\[
  (1 - t)^{-d} = \sum_{m \ge 0} \binom{m + d - 1}{d - 1} t^m.
\]
(This can of course be proved in other ways, for example by induction on $d$.)
If $f(t) = a_0 + a_1t + \cdots + a_st^s$ then
\[
  (*) \quad l(M_n) = a_0\binom{n+d-1}{d-1} + a_1\binom{n+d-2}{d-1} + \cdots + a_s\binom{n-s+d-1}{d-1},
\]
where we interpret $\binom{m}{d-1} = 0$ for $m < d - 1$. The right-hand side of $(*)$ can be
formally rearranged as a polynomial in $n$ with rational coefficients, say $\varphi(n)$;
then
\[
  \varphi(X) = \frac{f(1)}{(d-1)!}X^{d-1} + \text{(terms of lower degree)}.
\]
Since $\binom{m + d - 1}{d-1}$ coincides with the polynomial $m(m - 1) \cdots (m - d + 2)/
(d-1)!$ for $m \ge 0$, this implies the following result.

\begin{corollary*}
If $d_1 = \cdots = d_r = 1$ in Theorem 2, and $d = d(M)$ is defined
as above, then there is a polynomial $\varphi_M(X)$ of degree $d - 1$ with rational
coefficients such that for $n \ge s + 1 - d$ we have $l(M_n) = \varphi_M(n)$. Here $s$ is the
degree of the polynomial $(1 - t)^dP(M, t)$.
\end{corollary*}

The polynomial $\varphi_M$ appearing here is called the \emph{Hilbert polynomial}\index{Hilbert polynomial} of the
graded module $M$. The numerical function $l(M_n)$ itself is called the \emph{Hilbert}
\emph{function} of $M$; by the degree of a Hilbert function we mean the degree of the
corresponding Hilbert polynomial.

\begin{remark}
For general $d_1, \ldots, d_r$, it is no longer necessarily the case that $l(M_n)$
can be represented by one polynomial.
\end{remark}

\begin{example}\label{ex:13.1}
When $R = R_0[X_1, \ldots, X_r]$, the number of monomials of
degree $n$ is $\binom{n+r-1}{r-1}$, so that
\[
  l(R_n) = l(R_0)\binom{n+r-1}{r-1}
\]
holds for every $n \ge 0$, and the right-hand side is $\varphi_R(n)$. Thus $\varphi_R(X) =
(l(R_0)/r!)(X + r)(X + r - 1) \cdots (X + 1)$.
\end{example}

\begin{example}\label{ex:13.2}
Let $k$ be a field, and $F(X_1, \ldots, X_r)$ a homogeneous poly-
nomial of degree $s$; set $R = k[X_1, \ldots, X_r]/(F(X))$. Then for $n \ge s$,
\[
  l(R_n) = \binom{n+r-1}{r-1} - \binom{n-s+r-1}{r-1},
\]
and hence, setting $\binom{n+r-1}{r-1} = (1/r!)n^{r-1} + a_1n^{r-2} + \cdots$, we have
\[
  \varphi_R(X) = \frac{1}{r!}[X^{r-1} - (X - s)^{r-1}] + a_1[X^{r-2} - (X - s)^{r-2}] + \cdots
\]
\[
  = \frac{s}{(r-1)!}X^{r-2} + \text{(terms of lower degree)}.
\]
\end{example}

\begin{example}\label{ex:13.3}
Let $k$ be a field, and $R = k[X_1, \ldots, X_r]/P = k[\xi_1, \ldots, \xi_r]$,
where $P$ is a homogeneous prime ideal. Let $t$ be the transcendence degree
of $R$ over $k$, and suppose that $\xi_1, \ldots, \xi_t$ are algebraically independent
over $k$; then there are $\binom{n+t-1}{t-1}$ monomials of degree $n$ in the
$\xi_1, \ldots, \xi_t$, and these are linearly independent over $k$, so that
\[
  l(R_n) \ge \binom{n+t-1}{t-1},
\]
from which it follows that $d \ge t$. In fact we will prove
later (Theorem 8) that $d = t$.
\end{example}

A homogeneous ideal of the polynomial ring $k[X_1, \ldots, X_r]$ over a field
$k$ defines an algebraic variety in $r$-dimensional projective space $\mathbb{P}^r$, and
the Hilbert polynomial plays an important role in algebraic geometry.
For example, note that the numerator of the leading term of $\varphi_R$ in Example
2 is equal to the degree of $F$. This holds in more generality, but we must
leave details of this to textbooks on algebraic geometry [Ha].

The idea of using the construction of $\gr_m(A)$ to relate the study of a
general Noetherian local ring $(A, m)$ to the theory of ideals in a polynomial
ring over a field was one of the crucial ideas introduced by Krull in his
article `Dimension theory of local rings' [6], a work of monumental
significance for the theory of Noetherian rings. If $m$ is generated by $r$
elements then $\gr_m(A)$ is of the form $k[X_1, \ldots, X_r]/I$, where $k = A/m$ and $I$ is a
homogeneous ideal. However, the Hilbert function of this graded ring was
first used in the study of the multiplicity of $A$ by P. Samuel (1951).

\subsection*{Samuel functions}

In a little more generality, let $A$ be a Noetherian semilocal ring, and $m$
the Jacobson radical of $A$. If $I$ is an ideal of $A$ such that for some $\nu > 0$
we have $m^\nu \subset I \subset m$, we call $I$ an \emph{ideal of definition}\index{ideal of definition}; the $I$-adic and $m$-adic
topologies then coincide, so that `ideal of definition' means `ideal defining
the $m$-adic topology'. Let $M$ be a finite $A$-module. If we set
\[
  \gr_I(M) = \bigoplus_{n \ge 0} I^nM/I^{n+1}M
\]
then $\gr_I(M)$ is in a natural way a graded module over $\gr_I(A) = \bigoplus I^n/I^{n+1}$.
For brevity write $\gr_I(A) = A^*$ and $\gr_I(M) = M^*$. Then the ring $A^*_0 = A/I$ is
Artinian, and if $I = \sum_{i=1}^{r} x_iA$, and $\bar{x}_i$ is the image of $x_i$ in $I/I^2$, then
$A^* = A^*_0[\bar{x}_1, \ldots, \bar{x}_r]$. If also $M = \sum A\omega_i$ then $M^* = \sum A^*\bar{\omega}_i$ (where $\bar{\omega}_i$ is the
image of $\omega_i$ in $M_0^* = M/IM$), so that we can apply Theorem 2 and its
corollary to $M^*$. Noting that $l(M_n^*) = l(I^nM/I^{n+1}M)$ (where on the left-hand
side $l$ is the length as an $A^*_0$-module, on the right-hand side as an $A$-module),
we have
\[
  \sum_{n=0}^{\infty} l(M_n^*)t^n = \sum_{n=0}^{\infty} l(I^nM/I^{n+1}M)t^n.
\]
We now set $\chi_M^I(n) = l(M/I^{n+1}M)$. In particular we abbreviate $\chi_A^I(n)$ to
$\chi_A(n)$, and call it the \emph{Samuel function}\index{Samuel function} of the $A$-module $M$.

Repeatedly using the well-known formula
\[
  \binom{n}{r} = \binom{n-1}{r} + \binom{n-1}{r-1}
\]
we get
\[
  l(M/I^{n+1}M) = \sum_{k=0}^{n} l(I^kM/I^{k+1}M),
\]
so that from formula $(*)$ on p. 95 we get
\[
  \chi_M^I(n) = a_dn^d + a_{d-1}n^{d-1} + \cdots + a_0 \quad \text{with } a_i \in \Z.
\]
When $n \ge s$ this is a polynomial in $n$ of degree $d$. This degree $d$ is determined by $M$, and does not depend on $I$; to see this, if $I$ and $J$ are
both ideals of definition of $A$ then there exist natural numbers $a$ and $b$
such that $I^a \subset J$, $J^b \subset I$, so that
\[
  \chi_M^I(an + a - 1) \ge \chi_M^J(n) \quad \text{and} \quad \chi_M^J(bn + b - 1) \ge \chi_M^I(n).
\]
We thus write $d = d(M)$. It is natural to think of $d(M)$ as a measure of the size
of $M$.

\begin{theorem}\label{thm:13.3}
Let $A$ be a semilocal Noetherian ring, and $0 \to M' \to
M \to M'' \to 0$ an exact sequence of finite $A$-modules; then
\[
  d(M) = \max(d(M'), d(M'')).
\]
If $I$ is any ideal of definition of $A$, then $\chi_M^I$ and $\chi_{M''}^I$ have the same
leading coefficient.
\end{theorem}

\begin{proof}
We can assume $M'' = M/M'$. Then since $M''/I^nM'' = M/(M' + I^nM)$
we have
\[
  l(M/I^nM) = l(M/M' + I^nM) + l(M' + I^nM/I^nM)
\]
\[
  = l(M''/I^nM'') + l(M'/M' \cap I^nM).
\]
Thus setting $q(n) = l(M'/M' \cap I^{n+1}M)$, we have $\chi_M^I = \chi_{M''}^I + q$. Since
moreover both $\chi_{M''}^I$ and $q$ take on only positive values, $d(M)$ coincides with
whichever is the greater of $d(M'')$ and $\deg q$. However, by the Artin--Rees
lemma, there is a $c > 0$ such that
\[
  I^{n+1}M \cap M' \subset M' \cap I^{n+1}M \subset I^{n+1-c}M',
\]
and hence
\[
  \chi_{M''}^I(n) \ge q(n) \ge \chi_{M'}^I(n - c);
\]
therefore $q$ and $\chi_{M'}^I$ have the same leading coefficient.
\end{proof}

We now define a further measure $\delta(M)$ of the size of $M$: let $\delta(M)$ be the
smallest value of $n$ such that there exist $x_1, \ldots, x_n \in m$ for which $l(M/x_1M +
\cdots + x_nM) < \infty$. When $l(M) < \infty$, we interpret this as $\delta(M) = 0$. If $I$ is any
ideal of definition of $A$ then $l(M/IM) < \infty$, so that $\delta(M) \le$ number of
generators of $I$. Conversely, in the case that $A$ is a local ring and $M = A$,
then $l(A/I) < \infty$ implies that $I$ is an $m$-primary ideal. Therefore in this case
$\delta(A)$ is the minimum of the number of generators of $m$-primary ideals.

We have now arrived at the fundamental theorem of dimension theory.

\begin{theorem}\label{thm:13.4}
Let $A$ be a semilocal Noetherian ring and $M$ a finite
$A$-module; then we have
\[
  \dim M = d(M) = \delta(M).
\]
\end{theorem}

\begin{proof}
\textbf{Step 1.} Each of $d(M)$ and $\delta(M)$ are finite, but the finiteness of $\dim M$
has not yet been established. First of all, let us prove that $d(A) \ge \dim A$
for the case $M = A$, by induction on $d(A)$. Set $m = \rad(A)$. If $d(A) = 0$ then
$l(A/m^n)$ is constant for $n \gg 0$, so that for some $n$ we have $m^n = m^{n+1}$ and
by NAK, $m^n = 0$. Hence any prime ideal of $A$ is maximal and $\dim A = 0$.

Next suppose that $d(A) > 0$; if $\dim A = 0$ then we're done. If $\dim A > 0$,
consider a strictly increasing sequence $p_0 \subset p_1 \subset \cdots \subset p_e$ of prime ideals
of $A$, choose some element $x \in p_1 - p_0$, and set $B = A/(p_0 + xA)$; then by
the previous theorem applied to the exact sequence
\[
  0 \to A/p_0 \xrightarrow{x} A/p_0 \to B \to 0,
\]
we have $d(B) < d(A)$, and so by induction
\[
  \dim B \le d(B) \le d(A) - 1.
\]
(The values of $d(B)$ and $\dim B$ are independent of whether we consider $B$ as
an $A$-module or as a $B$-module, as is clear from the definitions.) In $B$, the
image of $p_1 \subset \cdots \subset p_e$ provides a chain of prime ideals of length $e - 1$, so that
\[
  e - 1 \le \dim B \le d(A) - 1;
\]
hence $e < d(A)$. Since this holds for any chain of prime ideals of $A$, this
proves $\dim A \le d(A)$. For general $M$, by Theorem~\ref{thm:6.4}, there are submodules
$M_i$ such that $0 = M_0 \subset M_1 \subset \cdots \subset M_r = M$ with
\[
  M_i/M_{i-1} \cong A/p_i, \quad \text{and} \quad p_i \in \Spec A.
\]
Since for an exact sequence $0 \to M' \to M \to M'' \to 0$ of finite $A$-
modules we have
\[
  \Supp(M) = \Supp(M') \cup \Supp(M'')
\]
and
\[
  \dim M = \max(\dim M', \dim M''),
\]
it is easy to see that
\[
  d(M) = \max\{d(A/p_i)\} \ge \max\{\dim(A/p_i)\} = \dim M.
\]

\textbf{Step 2.} We show that $\delta(M) \ge d(M)$. If $\delta(M) = 0$ then $l(M) < \infty$ so that $\chi_M(n)$
is bounded, hence $d(M) = 0$. Next suppose that $\delta(M) = s > 0$, choose
$x_1, \ldots, x_s \in m$ such that $l(M/x_1M + \cdots + x_sM) < \infty$, and set $M_i = M/x_1M +
\cdots + x_iM$; then clearly $\delta(M_i) = \delta(M) - i$. On the other hand,
\[
  l(M_i/m^\nu M_i) = l(M/x_iM + m^\nu M)
\]
\[
  = l(M/m^\nu M) - l(x_iM/x_iM \cap m^\nu M)
\]
\[
  = l(M/m^\nu M) - l(M/(m^\nu M:x_i))
\]
\[
  \ge l(M/m^\nu M) - l(M/m^{\nu-1}M),
\]
so that $d(M_i) \ge d(M) - 1$. Repeating this, we get $d(M_s) \ge d(M) - s$, but since
$\delta(M_s) = 0$ we have $d(M_s) = 0$, so that $s \ge d(M)$.

\textbf{Step 3.} We show that $\dim M \ge \delta(M)$, by induction on $\dim M$. If $\dim M = 0$
then $\Supp(M) \subset \text{m-Spec } A = V(m)$ so that for large enough $n$ we have
$m^n \subset \ann M$, and $l(M) < \infty$, therefore $\delta(M) = 0$. Next suppose that
$\dim M > 0$, and let $p_i$ for $1 \le i \le t$ be the minimal prime divisors of $\ann(M)$
with $\coht p_i = \dim M$; then the $p_i$ are not maximal ideals, so do not contain
$m$. Hence we can choose $x \in m$ not contained in any $p_i$. Setting
$M_1 = M/xM$ we get $\dim M_1 < \dim M$. Therefore by the inductive hypo-
thesis $\delta(M_1) \le \dim M_1$; but obviously $\delta(M) \le \delta(M_1) + 1$, so that $\delta(M) \le
\dim M_1 + 1 \le \dim M$.
\end{proof}

\begin{theorem}\label{thm:13.5}
Let $A$ be a Noetherian ring, and $I = (a_1, \ldots, a_r)$ an ideal
generated by $r$ elements; then if $p$ is a minimal prime divisor of $I$ we have
$\htop p \le r$. Hence the height of a proper ideal of $A$ is always finite.
\end{theorem}

\begin{proof}
The ideal $IA_p \subset A_p$ is a primary ideal belonging to the maximal
ideal, so that $\htop p = \dim A_p = \delta(A_p) \le r$.
\end{proof}

\begin{remark}
Krull proved this theorem by induction on $r$; the case $r = 1$ is
then the hardest part of the proof. Krull called the $r = 1$ case the
\emph{principal ideal theorem} (Hauptidealsatz), and the whole of Theorem 5 is sometimes
known by this name. Here Theorem 5 is merely a corollary of Theorem 4,
but one can also deduce the statement $\dim M = \delta(M)$ of Theorem 4 from
it. As far as proving Theorem 5 is concerned, Krull's proof, which does
not use the Samuel function, is easier. For this proof, see [N1] or [K].
More elementary proofs of the principal ideal theorem can be found in
Rees [3] and Caruth [1].

The definition of height is abstract, and even when one can find a lower
bound, one cannot expect an upper bound just from the definition, so
that this theorem is extremely important. The principal ideal theorem
corresponds to the familiar and obvious-looking proposition of geo-
metrical and physical intuition (which is strictly speaking not always true)
that `adding one equation can decrease the dimension of the space of
solutions by at most one'.
\end{remark}

\begin{theorem}\label{thm:13.6}
Let $P$ be a prime ideal of height $r$ in a Noetherian ring $A$.
Then
\begin{enumerate}[(i)]
  \item $P$ is a minimal prime divisor of some ideal $(a_1, \ldots, a_r)$ generated by
        $r$ elements;
  \item if $b_1, \ldots, b_s \in P$ we have $\htop P/(b_1, \ldots, b_s) \ge r - s$;
  \item if $a_1, \ldots, a_r$ are as in (i) we have
        \[
          \htop P/(a_1, \ldots, a_i) = r - i \quad \text{for } 1 \le i \le r.
        \]
\end{enumerate}
\end{theorem}

\begin{proof}
(i) $A_P$ is an $r$-dimensional local ring, so that by Theorem 4 we can
choose $r$ elements $a_1, \ldots, a_r \in PA_P$ such that $(a_1, \ldots, a_r)A_P$ is $PA_P$-primary.
Each $a_i$ is of the form an element of $P$ times a unit of $A_P$, so that without
loss of generality we can assume that $a_i \in P$. Then $P$ is a minimal prime
divisor of $(a_1, \ldots, a_r)A$.

(ii) Set $\bar{A} = A/(b_1, \ldots, b_s)$, $\bar{P} = P/(b_1, \ldots, b_s)$ and $\htop P = t$. Then by (i),
there exist $c_1, \ldots, c_t \in P$ such that $P$ is a minimal prime divisor of
$(c_1, \ldots, c_t)A$. Then $\bar{P}$ is a minimal prime divisor of $(\bar{b}_1, \ldots, \bar{b}_s, \bar{c}_1, \ldots, \bar{c}_t)$,
and hence $r \le s + t$ by Theorem 5.

(iii) In $\bar{A} = A/(a_1, \ldots, a_i)$, the ideal $\bar{P} = P/(a_1, \ldots, a_i)$ is a minimal
prime divisor of $(\bar{a}_{i+1}, \ldots, \bar{a}_r)$, hence
\[
  \htop P/(a_1, \ldots, a_i) \le r - i.
\]
The opposite inequality was proved in (ii).
\end{proof}

\begin{theorem}\label{thm:13.7}
Let $A = \bigoplus_{n \ge 0} A_n$ be a Noetherian graded ring.
\begin{enumerate}[(i)]
  \item If $I$ is a homogeneous ideal and $P$ is a prime divisor of $I$ then $P$ is also
        homogeneous.
  \item If $P$ is a homogeneous prime ideal of height $r$ then there exists a
        sequence
        \[
          P = P_r \supset P_{r-1} \supset \cdots \supset P_0
        \]
        of length $r$ consisting of homogeneous
        prime ideals.
\end{enumerate}
\end{theorem}

\begin{proof}
(i) $P$ can be expressed in the form $P = \ann(x)$ for a suitable element
$x$ of the graded $A$-module $A/I$. Let $a \in P$, and let $x = x_0 + x_1 + \cdots + x_n$
and $a = a_p + a_{p+1} + \cdots + a_q$ be decompositions into homogeneous terms.
Then since $ax = 0$,
\[
  a_px_0 = 0, \quad a_px_1 + a_{p+1}x_0 = 0, \quad a_px_2 + a_{p+1}x_1 + a_{p+2}x_0 = 0, \ldots,
\]
from which we get $a_p^2x_0 = 0$, $a_p^2x_1 = 0, \ldots$, and finally $a_p^{n+1}x = 0$. It follows
that $a_p^{n+1} \in P$, but since $P$ is prime, $a_p \in P$. Thus $a_{p+1} + \cdots + a_q \in P$, so
that in turn $a_{p+1} \in P$. Proceeding in the same way, we see that all the homo-
geneous terms of $a$ are in $P$, so that $P$ is a homogeneous ideal.

(ii) First of all note that we can assume that $A$ is an integral domain.
To see this, if we take a chain $P = P_r \supset \cdots \supset P_0$ of prime ideals of length
$r$ then $P_0$ is a minimal prime divisor of $(0)$, and so by (i) is a homogeneous
ideal; so we can replace $A$ by $A/P_0$. Now choose a homogeneous element
$0 \neq b_r \in P$; then by Theorem 6, $\htop (P/b_rA) = r - 1$, and so there is a minimal
prime divisor $Q$ of $b_rA$ such that $\htop (P/Q) = r - 1$; since $Q \neq (0)$, it is a
height $1$ homogeneous prime ideal. By the inductive hypothesis on $r$
applied to $P/Q$ there exists a chain $P = P_r \supset P_{r-1} \supset \cdots \supset P_1 = Q$ of homo-
geneous prime ideals of length $r - 1$, and adding on $P_0 = (0)$ we get a chain
of length $r$.
\end{proof}

Let us investigate more closely the relation between local rings and
graded rings.

\begin{theorem}\label{thm:13.8}
Let $k$ be a field, and $R = k[\xi_1, \ldots, \xi_r]$ a graded ring
generated by elements $\xi_1, \ldots, \xi_r$ of degree $1$; set $M = \sum \xi_iR$, $A = R_M$, and
$m = MA_M$.
\begin{enumerate}[(i)]
  \item Let $\chi$ be the Samuel function of the local ring $A$, and $\varphi$ the Hilbert
        function of the graded ring $R$; then $\varphi(n) = \chi(n) - \chi(n - 1)$;
  \item $\dim R = \htop M = \dim A = \deg \varphi + 1$;
  \item $\gr_m(A) \cong R$ as graded rings.
\end{enumerate}
\end{theorem}

\begin{proof}
$M$ is a maximal ideal of $R$ so that
\[
  m^n/m^{n+1} \cong M^n/M^{n+1} \cong R_n;
\]
hence $\chi(n) - \chi(n - 1) = l(m^n/m^{n+1}) = l(R_n) = \varphi(n)$, and so $\dim A = \deg \chi = 1 +
\deg \varphi$. Then since $A = R_M$, we have $\dim A = \htop M$. After this it is enough
to prove that $\dim R = \htop M$. First of all, assume that $R$ is an integral
domain, so that by Example 3 in the section on Hilbert functions and by
Theorem~\ref{thm:5.6}, we have
\[
  1 + \deg \varphi \ge \trdeg_k R = \dim R \ge \htop M;
\]
putting this together with $\htop M = \dim A = 1 + \deg \varphi$, we get $\dim R = \htop M$.
Next for general $R$, let $P_1, \ldots, P_s$ be the minimal prime ideals of $R$; then by
Theorem 7, these are all homogeneous ideals, and each $R/P_i$ is a graded
ring. Choosing $P_1$ such that $\dim R = \dim R/P_1$ and using the above result,
we get
\[
  \dim R = \dim R/P_1 = \htop M/P_1 \le \htop M \le \dim R,
\]
so that $\dim R = \htop M$ as required. We have $R_n \subset M^n \subset m^n$ with $m^n/m^{n+1} \cong
R_n$, and so taking an element $x$ of $R_n$ into its image in $m^n/m^{n+1}$ we obtain
a canonical one-to-one map $R \cong \gr_m A$, and it is clear from the definition
that this is a ring isomorphism.
\end{proof}

\begin{theorem}\label{thm:13.9}
Let $(A, m, k)$ be a Noetherian local ring, and set $G = \gr_m A$;
then $\dim A = \dim G$.
\end{theorem}

\begin{proof}
Letting $\varphi$ be the Hilbert polynomial of $G$, we have $\dim A =
1 + \deg \varphi$ (by Theorem 4), and by the previous theorem this is equal to
$\dim G$.
\end{proof}

In fact the following more general theorem holds: for $I$ a proper ideal in a
Noetherian local ring $A$, set $G = \gr_I(A)$; then $\dim A = \dim G$. This will be
proved a little later (Theorem~\ref{thm:15.7}).

% ============ 习题 ============
\section*{Exercises to \S\ref{sec:13}}
\emph{Prove the following propositions.}

\begin{exercise}\label{ex:13.1-ex}
Let $R = R_0 + R_1 + \cdots$ be a graded ring, and $u$ a unit of $R_0$. Then the map $T_u$,
defined by $T_u(x_0 + x_1 + \cdots + x_n) = x_0 + x_1u + \cdots + x_nu^n$ (where $x_i \in R_i$) is
an automorphism of $R$. If $R_0$ contains an infinite field $k$, then an ideal $I$ of $R$
is homogeneous if and only if $T_u(I) = I$ for every $u \in k$.
\end{exercise}

\begin{exercise}\label{ex:13.2-ex}
Let $R = R_0 + R_1 + \cdots$ be a graded ring, $I$ an ideal of $R$ and $t$ an
indeterminate over $R$. Set $R' = R[t, t^{-1}]$ and consider $R'$ as a graded ring
where $t$ has degree $0$ (that is, $R'_n = R_n[t, t^{-1}]$). Then an ideal $I$ of $R$ is
homogeneous if and only if $T_u(IR') = IR'$.
\end{exercise}

\begin{exercise}\label{ex:13.3-ex}
Let $A$ be a Noetherian ring having an embedded associated prime. If $a \in A$
is a non-zero divisor satisfying $\bigcap_{n \ge 1} a^nA = (0)$, then $A/(a)$ also has an
embedded associated prime.
\end{exercise}

\begin{exercise}\label{ex:13.4-ex}
Let $R = \bigoplus_{n \in \Z} R_n$ be a $\Z$-graded ring. For an ideal $I$ of $R$, let $I^*$ denote the
greatest homogeneous ideal of $R$ contained in $I$, that is the ideal of $R$
generated by all the homogeneous elements of $I$.
\begin{enumerate}[(i)]
  \item If $P$ is prime so is $P^*$.
  \item If $P$ is a homogeneous prime ideal and $Q$ is a $P$-primary ideal then $Q^*$
        is again $P$-primary.
\end{enumerate}
\end{exercise}

\begin{exercise}\label{ex:13.5-ex}
Let $R$ be a $\Z$-graded integral domain; write $S$ for the multiplicative set
consisting of all non-zero homogeneous elements of $R$. Then $R_S$ is a graded
ring, and its component of degree $0$ is a field $(R_S)_0 = K$; if $R \neq R_S$ then $R_S
\cong K[X, X^{-1}]$, where the degree of $X$ is the greatest common divisor of the
degrees of elements of $S$.
\end{exercise}

\begin{exercise}\label{ex:13.6-ex}
Let $R$ be an $\N$-graded ring and $P$ an inhomogeneous prime ideal of $R$; then
there are no prime ideals contained between $P^*$ and $P$. If $\htop P < \infty$ then
$\htop P = \htop P^* + 1$ (Matijevic--Roberts [1]).
\end{exercise}

% ============ Appendix to §13. Determinantal ideals (原书 pp.101-103) ============
\section*{Appendix to \S\ref{sec:13}. Determinantal ideals}
\phantomsection
\addcontentsline{toc}{section}{Appendix to \S\ref{sec:13}. Determinantal ideals}

Let $M = (a_{ij})$ be an $r \times s$ matrix $(r \le s)$ with elements $a_{ij}$ in a Noetherian ring
$A$, and let $I_t$ be the ideal of $A$ generated by the $t \times t$ minors (that is
$t$-determinants) of $M$. When $t = r$ and $A$ is a polynomial ring
$k[X_1, \ldots, X_s]$ over a field $k$, Macaulay proved that all the prime divisors of
$I_t$ have height $\le s - r + 1$ ([Mac], p. 54). In his Ph.D. thesis, Eagon
generalised this result as follows: for an arbitrary Noetherian ring $A$, every
minimal prime divisor of $I_t$ has height $\le (r - t + 1)(s - t + 1)$. The following
ingenious proof is taken from Eagon--Northcott [1]. We begin with some
preliminary observations.

The following operations on a matrix $M$ with elements in a ring $A$ are
called \emph{elementary row operations}: (1) permutation of the rows; (2) replacing
$C_i$ by $uC_i + vC_j$, where $C_i$ and $C_j$ $(i \neq j)$ are two distinct rows of $M$, $u$ is a unit
of $A$ and $v$ is an element of $A$; elementary column operations are defined
similarly. The ideal $I_t$ does not change under these operations. Now, if an
element of $M$ is a unit in $A$, we can transform $M$ by a finite number of
elementary row and column operations to the following form:
\[
  \begin{pmatrix}
    1 & 0 & \cdots & 0 \\
    0 & & & \\
    \vdots & & N & \\
    0 & & &
  \end{pmatrix},
\]
and $I_t$ is equal to the ideal of $A$ generated by the $(t - 1) \times (t - 1)$ minors of $N$.

\begin{lemma}\label{lem:13.1}
Let $(A, P)$ be a Noetherian local ring and set $B = A[X]$. Let $J$ be a
$P$-primary ideal of $A$ and $J'$ an ideal of $B$ such that $J' \subset PB$ and
$J' + XB = JB + XB$. Then $PB$ is a minimal prime divisor of $J'$.
\end{lemma}

\begin{proof}
$PB + XB$ is a maximal ideal of $B$, and is the radical of $JB + XB =
J' + XB$. Thus, in the ring $B/J'$ we have that $(PB + XB)/J'$ is a minimal
prime divisor of the principal ideal $(J' + XB)/J'$. Hence $\htop((PB + XB)/
J') = 1$. Since $PB \subset J'$ is a prime ideal in $B/J'$, we have $\htop(PB/J') = 0$.
\end{proof}

\begin{theorem}[Eagon]\label{thm:13.10}
Let $A$ be a Noetherian ring and $M$ be an $r \times s$
matrix $(r \le s)$ of elements of $A$. Let $I_t$ be the ideal of $A$ generated by the $t \times t$
minors of $M$. If $P$ is a minimal prime divisor of $I_t$, then we have
\[
  \htop P \le (r - t + 1)(s - t + 1).
\]
\end{theorem}

\begin{proof}
Induction on $r$. When $r = 1$ we have $t = 1$, and so $(r - t + 1)(s - t +
1) = s$. The ideal $I_t$ is generated by $s$ elements, so that the assertion is just
the principal ideal theorem (Theorem 5) in this case. Next assume that $r > 1$.
Localising at $P$ we may assume that $A$ is a local ring with maximal ideal $P$,
and that $I_t$ is $P$-primary.

If $t = 1$, then $I_t$ is generated by $rs$ elements and $(r - t + 1)(s - t + 1) = rs$,
so our assertion holds also for this case. Therefore we assume $t > 1$. If at
least one of the elements of $M$ is a unit of $A$, then by what we said above, $I_t$ is
generated by $(t - 1) \times (t - 1)$ minors of a $(r - 1) \times (s - 1)$ matrix, and again
we are done. Therefore we assume that all the elements of $M$ are in $P$. Now
comes the brilliant idea. Let $M'$ be the matrix with elements in $B = A[X]$
obtained from $M$ by replacing $a_{11}$ by $a_{11} + X$, and let $I'$ be the ideal of $B$
generated by the $t \times t$ minors of $M'$. Since $t > 1$ and $a_{ij} \in P$ for all $i$ and $j$ we
have $I' \subset PB$. We also have $I' + XB = I_tB + XB$ since both sides have the
same image in $B/XB = A$. Therefore $PB$ is a minimal prime divisor of $I'$ by
the lemma. Since the element $a_{11} + X$ of $M'$ is not in $PB$, we have
$\htop PB \le (r - t + 1)(s - t + 1)$ by our previous argument. Since $\htop PB = \htop P$,
as we can see by Theorems 4 and 5, we are done.
\end{proof}

% ============ §14 Systems of parameters and multiplicity (原书 pp.104-110) ============
\section{Systems of parameters and multiplicity}\label{sec:14}\index{systems of parameters, section 14}\index{multiplicity, section 14}

Let $(A, m)$ be an $r$-dimensional Noetherian local ring; by Theorem~\ref{thm:13.4}, there exists an $m$-primary ideal generated by $r$ elements, but none
generated by fewer. If $a_1, \ldots, a_r \in m$ generate an $m$-primary ideal,
$\{a_1, \ldots, a_r\}$ is said to be a \emph{system of parameters}\index{system of parameters} of $A$ (sometimes
abbreviated to s.o.p.). If $M$ is a finite $A$-module with $\dim M = s$, there exist
$y_1, \ldots, y_s \in m$ such that $l(M/(y_1, \ldots, y_s)M) < \infty$, and then $\{y_1, \ldots, y_s\}$ is
said to be a \emph{system of parameters of $M$}.

If we set $A/m = k$, the smallest number of elements needed to generate
$m$ itself is equal to $\rank_k m/m^2$; (here $\rank_k$ is the rank of a free module
over $k$, that is the dimension of $m/m^2$ as $k$-vector space). This number is
called the \emph{embedding dimension}\index{embedding dimension} of $A$, and is written emb dim $A$. In general
\[
  \dim A \le \text{emb dim } A,
\]
and equality holds when $m$ can be generated by $r$ elements; in this case
$A$ is said to be a \emph{regular local ring}\index{regular local ring}, and a system of parameters generating
$m$ is called a \emph{regular system of parameters}.

\begin{theorem}\label{thm:14.1}
Let $(A, m)$ be a Noetherian local ring, and $x_1, \ldots, x_r$ a system
of parameters. Then
\begin{enumerate}[(i)]
  \item $\dim A/(x_1, \ldots, x_i) = r - i$ for $1 \le i \le r$.
  \item although it is not true that $\htop (x_1, \ldots, x_i) = i$ for all $i$ for an arbitrary
        system of parameters, there exists a choice of $x_1, \ldots, x_r$ such that every
        subset $F \subset \{x_1, \ldots, x_r\}$ generates an ideal of $A$ of height equal to the
        number of elements of $F$.
\end{enumerate}
\end{theorem}

\begin{proof}
(i) is contained in Theorem~\ref{thm:13.6}. We now prove the second half of
(ii). If $r \le 1$ the assertion is obvious; suppose that $r > 1$. Let $p_{0j}$ (for $1 \le
j \le e_0$) be the prime ideals of $A$ of height $0$. Choosing $x_1 \in m$ not contained
in any $p_{0j}$, we have $\htop (x_1) = 1$. Next letting $p_{1j}$ (for $1 \le j \le e_1$) be the minimal
prime divisors of $(x_1)$, so that $\htop p_{1j} = 1$, and choosing $x_2 \in m$ not contained
in any $p_{0j}$ or any $p_{1j}$, we have $\htop(x_1) = 1$, $\htop(x_1, x_2) = 2$; if $r = 2$ we're
done. If $r > 2$ we choose $x_3 \in m$ not contained in any minimal prime
divisor of $(0)$, $(x_1)$, $(x_2)$, $(x_1, x_2)$, and proceed in the same way to obtain
the result.

We now give an example where $\htop (x_1, \ldots, x_i) < i$. Let $k$ be a field and
set $R = k[X, Y, Z]$; let $I = (X) \cap (Y, Z)$, and write $A = R/I$, and $x, y, z$ for
the images in $A$ of $X, Y, Z$. The minimal prime ideals of $A$ are $(x)$ and
$(y, z)$; now $A/(x) = R/(X) \cong k[Y, Z]$ is two-dimensional and $A/(y, z) \cong
R/(Y, Z) \cong k[X]$ is one-dimensional, so that $\dim A = 2$. $\{y, x + z\}$ is a
system of parameters of $A$; in fact $xy = xz = 0$, so that $x^2 = x(x + z) \in
(y, x + z)$ and $z^2 = z(x + z) \in (y, x + z)$. However, $y$ is contained in the
minimal prime ideal $(y, z)$ of $A$, and hence $\htop (y) = 0$.
\end{proof}

\begin{theorem}\label{thm:14.2}
Let $(R, m)$ be an $n$-dimensional regular local ring, and
$x_1, \ldots, x_i$ elements of $m$. Then the following conditions are equivalent:
\begin{enumerate}[(1)]
  \item $x_1, \ldots, x_i$ is a subset of a regular system of parameters of $R$;
  \item the images in $m/m^2$ of $x_1, \ldots, x_i$ are linearly independent over $R/m$;
  \item $R/(x_1, \ldots, x_i)$ is an $(n - i)$-dimensional regular local ring.
\end{enumerate}
\end{theorem}

\begin{proof}
$(1) \Rightarrow (2)$ If $x_1, \ldots, x_i, x_{i+1}, \ldots, x_n$ is a regular system of parameters then
their images generate $m/m^2$ over $k = R/m$, and since $\rank_k m/m^2 = n$ they
must be linearly independent over $k$.

$(1) \Rightarrow (3)$ We know that $\dim R/(x_1, \ldots, x_i) = n - i$, and the images of
$x_{i+1}, \ldots, x_n$ generate the maximal ideal of $R/(x_1, \ldots, x_i)$.

$(3) \Rightarrow (1)$ If the maximal ideal $m/(x_1, \ldots, x_i)$ of $R/(x_1, \ldots, x_i)$ is generated
by the images of $y_1, \ldots, y_{n-i} \in m$ then $m$ is generated by $x_1, \ldots, x_i$,
$y_1, \ldots, y_{n-i}$.
\end{proof}

\begin{remark}
The hypothesis that $R$ is regular is not needed for $(3) \Rightarrow (1)$.
\end{remark}

\begin{proof}[Proof of $(2) \Rightarrow (1)$]
Using $\rank_k m/m^2 = n$, if we choose $x_{i+1}, \ldots, x_n \in m$ such that
the images of $x_1, \ldots, x_n$ in $m/m^2$ form a basis then $x_1, \ldots, x_n$ generate $m$
and so forms a regular system of parameters.
\end{proof}

\begin{theorem}\label{thm:14.3}
A regular local ring is an integral domain.
\end{theorem}

\begin{proof}
Let $(R, m)$ be an $n$-dimensional regular local ring; we proceed by
induction on $n$. If $n = 0$ then $m$ is an ideal generated by $0$ elements, so
that $m = (0)$. This in turn means that $R$ is a field. Thus a zero-dimensional
regular local ring is just a field by another name.

When $n = 1$, the maximal ideal $m = xR$ is principal and $\htop m = 1$, so that
there exists a prime ideal $p \neq m$ with $m \supset p$. If $y \in p$ we can write $y = xa$ with
$a \in R$, and since $x \notin p$ we have $a \in p$; hence $p = xp$, and by NAK, $p = (0)$. This
proves that $R$ is an integral domain. (There is a slightly different proof in the
course of the proof of Theorem~\ref{thm:11.2}; as proved there, a one-dimensional
regular local ring is just a DVR by another name.)

When $n > 1$, let $p_1, \ldots, p_l$ be the minimal prime ideals of $R$; then since
$m \nsubseteq m^2$ and $m \nsubseteq p_i$ for all $i$, there exists an element $x \in m$ not contained
in any of $m^2, p_1, \ldots, p_l$ (see Ex.~\ref{ex:1.6-ex}). Then the image of $x$ in $m/m^2$ is non-zero,
so that by the previous theorem $R/xR$ is an $(n - 1)$-dimensional regular
local ring. By the induction hypothesis, $R/xR$ is an integral domain, in
other words, $xR$ is a prime ideal of $R$. If $p_1$ is one of the minimal prime ideals
contained in $xR$ then since $x \notin p_1$, the same argument as in the $n = 1$ case
shows that $p_1 = xp_1$, and hence $p_1 = (0)$.
\end{proof}

\begin{theorem}\label{thm:14.4}
Let $(A, m, k)$ be a $d$-dimensional regular local ring; then
\[
  \gr_m(A) \cong k[X_1, \ldots, X_d],
\]
and if $\chi(n)$ is the Samuel function of $A$ then
\[
  \chi(n) = \binom{n+d}{d} \quad \text{for all } n \ge 0.
\]
\end{theorem}

\begin{proof}
Since $m$ is generated by $d$ elements, $\gr_m(A)$ is of the form
$k[X_1, \ldots, X_d]/I$, where $I$ is a homogeneous ideal. Now if $I \neq (0)$ let $S \in I$ be a
non-zero homogeneous element of degree $r$; then for $n \ge r$ the homogeneous
piece of $k[X]/I$ of degree $n$ has length at most
\[
  \binom{n+r-1}{d-1} - \binom{n-1}{d-1},
\]
which is a polynomial of degree $d - 2$ in $n$. This
implies that the Samuel function of $A$ is of degree at most $d - 1$, and
contradicts $\dim A = d$. Hence $I = (0)$; the second assertion follows from
the first.
\end{proof}

Let $(A, m)$ be a Noetherian local ring. Elements $y_1, \ldots, y_l \in m$ are said to be
\emph{analytically independent}\index{analytically independent elements} if they have the following property; for every
homogeneous form $F(Y_1, \ldots, Y_l)$ with coefficients in $A$,
\[
  F(y_1, \ldots, y_l) = 0 \Rightarrow \text{the coefficients of } F \text{ are in } m.
\]
If $y_1, \ldots, y_l$ are analytically independent and $A$ contains a field $k$, then
$F(y) \neq 0$ for any non-zero homogeneous form $F(Y) \in k[Y_1, \ldots, Y_l]$.

\begin{theorem}\label{thm:14.5}
Let $(A, m)$ be a $d$-dimensional Noetherian local ring and
$x_1, \ldots, x_d$ a system of parameters of $A$; then $x_1, \ldots, x_d$ are analytically
independent.
\end{theorem}

\begin{proof}
Set $q = \sum x_iA$. Since $q$ is an ideal of definition of $A$, by
Theorem~\ref{thm:13.4}, $\chi_q(n) = l(A/q^n)$ is a polynomial of degree $d$ in $n$ for $n \gg 0$. Set
$\bar{A} = k$; we say that a homogeneous form $\bar{f}(X) \in k[X_1, \ldots, X_d]$ of degree $n$
is a \emph{null-form} of $q$ if $F(x_1, \ldots, x_d) \in q^{n+1}m$ for any homogeneous form
$F(X) \in A[X_1, \ldots, X_d]$ which reduces to $\bar{f}(X)$ modulo $m$. Write $\mathfrak{n}$ for the ideal
of $k[X_1, \ldots, X_d]$ generated by the null-forms of $q$. Then
\[
  k[X_1, \ldots, X_d]/\mathfrak{n} = \bigoplus_{n \ge 0} q^n/q^{n+1}m,
\]
and writing $\varphi$ for the Hilbert polynomial of $k[X]/\mathfrak{n}$, we have $\varphi(n)
= l(q^n/q^{n+1}m)$ for $n \gg 0$. The right-hand side is just the number of elements in a
minimal basis of $q^n$, so that $\varphi(n) \cdot l(A/q) \ge l(q^n/q^{n+1})$. Now
\[
  l(q^n/q^{n+1}) = \chi_q(n) - \chi_q(n - 1)
\]
is a polynomial in $n$ of degree $d - 1$, so that $\deg \varphi \ge d - 1$, but if $\mathfrak{n} \neq (0)$ this
is impossible. Thus $\mathfrak{n} = (0)$, and the statement in the theorem follows at
once.
\end{proof}

\subsection*{Multiplicity}

Let $(A, m)$ be a $d$-dimensional Noetherian local ring, $M$ a finite $A$-module,
and $q$ an ideal of definition of $A$ (that is, an $m$-primary ideal). As we saw
in \S\ref{sec:13}, the Samuel function $l(M/q^{n+1}M) = \chi_M^q(n)$ can be expressed for
$n \gg 0$ as a polynomial in $n$ with rational coefficients, and degree equal to
$\dim M$, and therefore at most $d$. In addition, this polynomial can only take
integer values for $n \gg 0$, so it is easy to see by induction on $d$ (using the fact
that $\chi(n + 1) - \chi(n)$ has the same property) that
\[
  \chi_M^q(n) = \frac{e}{d!}n^d + \text{(terms of lower order)},
\]
with $e \in \Z$. This integer will be written $e(q, M)$. By definition we have the
following property.

\smallskip
\noindent\textbf{Formula 14.1.}\quad $\displaystyle e(q, M) = \lim_{n \to \infty} \frac{d!}{n^d}l(M/q^nM)$, and in particular, if $d = 0$ then
$e(q, M) = l(M)$.

\smallskip
\noindent From this we see easily the following:
\noindent\textbf{Formula 14.2.}\quad $e(q, M) > 0$ if $\dim M = d$, and $e(q, M) = 0$ if $\dim M < d$;

\smallskip
\noindent\textbf{Formula 14.3.}\quad $e(q^r, M) = e(q, M)r^d$;

\smallskip
\noindent\textbf{Formula 14.4.}\quad If $q$ and $q'$ are both $m$-primary ideals and $q \supset q'$ then
$e(q, M) \le e(q', M)$.

\smallskip
\noindent We set $e(q, A) = e(q)$, and define this to be the \emph{multiplicity} of $q$. In addition,
we will refer to the multiplicity $e(m)$ of the maximal ideal as the \emph{multiplicity of the
local ring} $A$, and sometimes write $e(A)$ for it. For example, if $A$ is a
regular local ring then by Theorem 4, we can see that $e(A) = 1$.

\begin{theorem}\label{thm:14.6}
Let $0 \to M' \to M \to M'' \to 0$ be an exact sequence of
finite $A$-modules. Then
\[
  e(q, M) = e(q, M') + e(q, M'').
\]
\end{theorem}

\begin{proof}
We view $M'$ as a submodule of $M$. Then
\[
  l(M/q^nM) = l(M''/q^nM'') + l(M'/M' \cap q^nM),
\]
and obviously $q^nM' \subset M' \cap q^nM$. On the other hand by Artin--Rees, there
exists $c > 0$ such that
\[
  M' \cap q^nM \subset q^{n-c}M' \quad \text{for all } n > c.
\]
Hence
\[
  l(M'/q^{n-c}M') \le l(M'/M' \cap q^nM) \le l(M'/q^nM').
\]
From this and Formula 14.1 it follows easily that
\[
  e(q, M) - e(q, M'') = \lim_{n \to \infty} \frac{d!}{n^d}l(M'/M' \cap q^nM) = e(q, M'). \qquad\blacksquare
\]
\end{proof}

\begin{theorem}\label{thm:14.7}
Let $\{p_1, \ldots, p_t\}$ be all the minimal prime ideals of $A$ such
that $\dim A/p_i = d$; then
\[
  e(q, M) = \sum_{i=1}^{t} l(M_{p_i})e(\bar{q}, A/p_i),
\]
where $\bar{q}_i$ denotes the image of $q$ in $A/p_i$, and $l(M_{p_i})$ stands for the length of $M_{p_i}$ as
$A_{p_i}$-module.
\end{theorem}

\begin{proof}[Proof (taken from Nagata [N1])]
We write $\sigma = \sum l(M_{p_i})$ and proceed by
induction on $\sigma$. If $\sigma = 0$ then $\dim M < d$, so that the left-hand side is $0$, and
the right-hand side is obviously $0$; now suppose $\sigma > 0$. Now there is some
$p \in \{p_1, \ldots, p_t\}$ for which $M_p \neq 0$; then $p$ is a minimal element of
$\Supp(M)$. Hence $p \in \Ass(M)$, that is $M$ contains a submodule $N$ isomorphic
to $A/p$. Then
\[
  e(q, M) = e(q, N) + e(q, M/N).
\]
On the other hand, $N_{p_i} \cong A_p/pA_p$, and $N_{p_i} = 0$ for $p_i \neq p$, so that $l(N_{p_i}) = 1$,
and the value of $\sigma$ for $M/N$ has decreased by one, so that the theorem holds
for $M/N$. However, from the definition
\[
  e(q, N) = e(q, A/p) = e(\bar{q}, A/p), \quad \text{where } \bar{q} = (q + p)/p.
\]
Putting this together, we see that the theorem also holds for $M$.
\end{proof}

Theorem 7 allows us to reduce the study of $e(q, M)$ to the case that $A$ is an
integral domain and $M = A$. In particular, if $A$ is an integral domain then
$l(M_{p_i})$ is just the rank of $M$, so that we obtain the following theorem.

\begin{theorem}\label{thm:14.8}
Let $A$ be a Noetherian local integral domain, $q$ an ideal of
definition of $A$ and $M$ a finite $A$-module; then
\[
  e(q, M) = e(q) \cdot s, \quad \text{where } s = \rank M.
\]
\end{theorem}

\begin{theorem}\label{thm:14.9}
Let $(A, m)$ be a Noetherian local ring, $q$ an ideal of definition
of $A$, and $x_1, \ldots, x_d$ a system of parameters of $A$ contained in $q$. Suppose that
$x_i \in q^{\nu_i}$ for $1 \le i \le d$. Then for a finite $A$-module $M$ and $s = 1, \ldots, d$ we have
\[
  e(q/(x_1, \ldots, x_s), M/(x_1, \ldots, x_s)M) \ge \nu_1\nu_2 \cdots \nu_s e(q, M).
\]
In particular if $s = d$, we have
\[
  l(M/(x_1, \ldots, x_d)M) \ge \nu_1\nu_2 \cdots \nu_de(q, M).
\]
\end{theorem}

\begin{proof}
It is enough to prove the case $s = 1$. We set $A' = A/x_1A$, $q' = q/x_1A$,
$M' = M/x_1M$ and $\nu = \nu_1$. By Theorem 1, we have $\dim A' = d - 1$. On the
other hand,
\[
  l(M'/q'^nM') = l(M/x_1M + q^nM)
\]
\[
  = l(M/q^nM) - l(x_1M + q^nM/q^nM).
\]
In addition, in view of $(x_1M + q^nM)/q^nM \cong x_1M/x_1M \cap q^nM \cong M/
(q^nM:x_1)$ and $q^{n-\nu}M \subset q^nM:x_1$, we have
\[
  l(M'/q'^nM') \ge l(M/q^nM) - l(M/q^{n-\nu}M),
\]
and therefore
\[
  l(M'/q'^nM') \ge l(M/q^nM) - l(M/q^{n-\nu}M).
\]
When $n \gg 0$ the right-hand side is of the form
\[
  \frac{e(q, M)}{d!}[n^d - (n - \nu)^d] + \text{(polynomial of degree } d - 2 \text{ in } n)
\]
\[
  = \frac{e(q, M)\nu}{(d-1)!}n^{d-1} + \text{(polynomial of degree } d - 2 \text{ in } n),
\]
so that the assertion is clear.
\end{proof}

A case of the above theorem which is particularly simple, but important,
is the following.

\begin{theorem}\label{thm:14.10}
Let $(A, m)$ be a $d$-dimensional Noetherian local ring, let
$x_1, \ldots, x_d$ be a system of parameters of $A$, and set $q = (x_1, \ldots, x_d)$; then
\[
  l(A/q) \ge e(q) \ge e(m).
\]
and if in addition $x_i \in m^{\nu_i}$ for all $i$ then $l(A/q) \ge \nu_1 \cdots \nu_d e(m)$.
\end{theorem}

\begin{theorem}\label{thm:14.11}
Let $A, m, x_i$ and $q$ be as above. Let $M$ be a finite $A$-module,
and set $A' = A/x_1A$, $M' = M/x_1M$ and $q' = q/x_1A = \sum_{i \ge 2} x_iA'$. Then if $x_1$ is a
non-zero-divisor of $M$, we have the following equality
\[
  e(q', M') = e(q, M).
\]
\end{theorem}

\begin{proof}
Since $l(M'/q'^{n+1}M') = l(M/x_1M + q^{n+1}M)$ we have
\[
  l(M/q^{n+1}M) - l(M'/q'^{n+1}M') = l(x_1M + q^{n+1}M/q^{n+1}M)
\]
\[
  = l(x_1M/x_1M \cap q^{n+1}M) = l(M/(q^{n+1}M:x_1))
\]
\[
  = l(M/q^nM) - l((q^{n+1}M:x_1)/q^nM).
\]
On the other hand, setting $\mathfrak{a} = \sum_{i \ge 2} x_iA$ we have $q = x_1A + \mathfrak{a}$ and $q^{n+1} =
x_1q^n + \mathfrak{a}^{n+1}$, and therefore
\[
  q^{n+1}M:x_1 = q^nM + (\mathfrak{a}^{n+1}M:x_1).
\]
Moreover, by Artin--Rees, there is a $c > 0$ such that for $n > c$ we have
$\mathfrak{a}^{n+1}M \cap x_1M = \mathfrak{a}^{n-c}(\mathfrak{a}^{c+1}M \cap x_1M)$, and therefore $\mathfrak{a}^{n+1}M:x_1 \subset \mathfrak{a}^{n-c}M$.
Thus
\[
  (q^{n+1}M:x_1)/q^nM = (q^nM + (\mathfrak{a}^{n+1}M:x_1))/q^nM
\]
\[
  \subset (q^nM + \mathfrak{a}^{n-c}M)/q^nM
\]
\[
  \cong \mathfrak{a}^{n-c}M/\mathfrak{a}^{n-c}M \cap q^nM.
\]
Now $\mathfrak{a}^{n-c}M/\mathfrak{a}^{n-c}M \cap q^nM$ is a module over $A/q'$, and since $\mathfrak{a}$ is generated
by $d - 1$ elements, $\mathfrak{a}^{n-c}$ is generated by $\binom{n-c+d-2}{d-2}$ elements. Thus
for $n > c$ we have
\[
  l(M'/q'^{n+1}M') \ge l(M/q^{n+1}M) - m\binom{n-c+d-2}{d-2},
\]
where $m$ is the number of generators of $M$. The right-hand side is a
polynomial of degree $d - 2$ in $n$, so that
\[
  e(q', M') = (d - 1)! \lim_{n \to \infty} l(M'/q'^{n+1}M')/n^{d-1}
\]
\[
  = (d - 1)! \lim_{n \to \infty} [l(M/q^{n+1}M) - l(M/q^nM)]/n^{d-1}
\]
\[
  = e(q, M).
\]
\end{proof}

\begin{theorem}[Lech's lemma]\label{thm:14.12}
Let $A$ be a $d$-dimensional Noetherian
local ring, and $x_1, \ldots, x_d$ a system of parameters of $A$; set $q = (x_1, \ldots, x_d)$, and
suppose that $M$ is a finite $A$-module. Then
\[
  e(q, M) = \lim_{\min(\nu_i) \to \infty} \frac{l(M/(x_1^{\nu_1}, \ldots, x_d^{\nu_d})M)}{\nu_1 \cdots \nu_d}.
\]
\end{theorem}

\begin{proof}
If $d = 0$ then both sides are equal to $l(M)$. If $d = 1$ then the right-hand
side is exactly Formula 14.1 which defines $e(q, M)$. For $d > 1$ we use
induction on $d$.

Setting $N_\nu = \{m \in M \mid x_1^\nu m = 0\}$ we have $N_1 \subset N_2 \subset \cdots$ so that there is a
$c > 0$ such that $N_\nu = N_{\nu+1} = \cdots$. If we set $M' = x_1^cM$ then $x_1$ is a non-zero-
divisor for $M'$, and there is an exact sequence $0 \to N_\nu \to M \to M' \to 0$.
Since $N_\nu$ is a module over $A/x_1^cA$ we have $\dim N_\nu < d$, and therefore
$e(q, N_\nu) = e(q, M')$. On the other hand,
\[
  l(M/(x_1^{\nu_1}, \ldots, x_d^{\nu_d})M) - l(M'/(x_1^{\nu_1}, \ldots, x_d^{\nu_d})M')
\]
\[
  = l(N_\nu + (x_1^{\nu_1}, \ldots, x_d^{\nu_d})M/(x_1^{\nu_1}, \ldots, x_d^{\nu_d})M)
\]
\[
  = l(N_\nu/N_\nu \cap (x_1^{\nu_1}, \ldots, x_d^{\nu_d})M)
\]
\[
  \le l(N_\nu/(x_1^{\nu_1}, \ldots, x_d^{\nu_d})N_\nu).
\]
If $\nu_1 \ge c$ then $x_1^{\nu_1}N_\nu = 0$, and $N_\nu$ is a module over the $(d - 1)$-dimensional
local ring $A/x_1^cA$, so that by induction there is a constant $C$ such that as
$\min(\nu_i) \to \infty$ we have
\[
  l(N_\nu/(x_1^{\nu_1}, \ldots, x_d^{\nu_d})N_\nu) = l(N_\nu/(x_2^{\nu_2}, \ldots, x_d^{\nu_d})N_\nu) \le C\nu_2 \cdots \nu_d.
\]
Therefore,
\[
  \lim [l(M/(x_1^{\nu_1}, \ldots, x_d^{\nu_d})M) - l(M'/(x_1^{\nu_1}, \ldots, x_d^{\nu_d})M')]/\nu_1 \cdots \nu_d = 0.
\]
This means that we can replace $M$ by $M'$ in the theorem, and so we can
assume that $x_1$ is a non-zero-divisor in $M$. Then by the previous theorem we
have $e(q, M) = e(\bar{q}, \bar{M})$, with $\bar{q} = q/x_1A$ and $\bar{M} = M/x_1M$. If we furthermore set
\[
  E = (x_2^{\nu_2}, \ldots, x_d^{\nu_d})M \quad \text{and} \quad F = M/E
\]
then by Theorem 9, we have
\[
  e(\bar{q}, \bar{M}) \cdot \nu_2 \cdots \nu_d \le l(M/(x_1^{\nu_1}, \ldots, x_d^{\nu_d})M) = l(F/x_1^{\nu_1}F)
\]
\[
  = \sum_{i=1}^{\nu_1} l(x_1^{i-1}F/x_1^iF) \le \nu_1l(F/x_1F) = \nu_1l(M/(x_1, x_2^{\nu_2}, \ldots, x_d^{\nu_d})M).
\]
Then by induction on $d$ we have
\[
  \lim l(M/(x_1^{\nu_1}, \ldots, x_d^{\nu_d})M)/\nu_1 \cdots \nu_d = \lim l(\bar{M}/(\bar{x}_2^{\nu_2}, \ldots, \bar{x}_d^{\nu_d})\bar{M})/\nu_2 \cdots \nu_d
  = e(\bar{q}, \bar{M}). \qquad\blacksquare
\]
\end{proof}

Although we will not use it in this book, we state here without proof a
remarkable result of Serre which shows that multiplicity can be expressed
as the Euler characteristic of the homology groups of the Koszul complex
(discussed in \S\ref{sec:16}).

\begin{theorem*}
Let $A$ be a $d$-dimensional Noetherian local ring, and $x_1, \ldots, x_d$ a
system of parameters of $A$; set $q = (x_1, \ldots, x_d)$ and let $M$ be a finite $A$-
module. Then
\[
  e(q, M) = \sum_i (-1)^i l(H_i(x, M)).
\]
For a proof, see for example Auslander and Buchsbaum [2].
\end{theorem*}

As we have seen in several of the above theorems, the multiplicity of
ideals generated by systems of parameters enjoy various nice properties. We
are now going to see that in a certain sense the general case can be reduced
to this one. We follow the method of Northcott and Rees [1].

More generally, let $A$ be a ring and $\mathfrak{a}$ an ideal. We say that an ideal $\mathfrak{b}$ is a
\emph{reduction}\index{reduction of an ideal} of $\mathfrak{a}$ if it satisfies the following condition:
\[
  \mathfrak{b} \subset \mathfrak{a}, \quad \text{and for some } r \ge 0 \text{ we have } \mathfrak{a}^{r+1} = \mathfrak{b}\mathfrak{a}^r.
\]
If $\mathfrak{b}$ is a reduction of $\mathfrak{a}$ and $\mathfrak{a}^{r+1} = \mathfrak{b}\mathfrak{a}^r$ then for any $n \ge 0$ we have $\mathfrak{a}^{r+n} = \mathfrak{b}^n\mathfrak{a}^r$.

\begin{theorem}\label{thm:14.13}
Let $(A, m)$ be a Noetherian local ring, $q$ an $m$-primary ideal
and $\mathfrak{b}$ a reduction of $q$; then $\mathfrak{b}$ is also $m$-primary, and for any finite $A$-module
$M$ we have
\[
  e(q, M) = e(\mathfrak{b}, M).
\]
\end{theorem}

\begin{proof}
If $q^{r+1} = \mathfrak{b}q^r$ then $q^{r+1} \subset \mathfrak{b} \subset q$, hence $\mathfrak{b}$ is also $m$-primary.
Moreover,
\[
  l(M/\mathfrak{b}^{n+1}M) \ge l(M/q^{n+1}M) = l(M/\mathfrak{b}^nq^r) \ge l(M/\mathfrak{b}^{n+r}M),
\]
so that $e(q, M) = e(\mathfrak{b}, M)$ follows easily.
\end{proof}

\begin{theorem}\label{thm:14.14}
Let $(A, m)$ be a $d$-dimensional Noetherian local ring, and
suppose that $A/m$ is an infinite field; let $q = (u_1, \ldots, u_s)$ be an $m$-primary
ideal. Then if $y_i = \sum a_{ij}u_j$ for $1 \le i \le d$ are $d$ `sufficiently general' linear
combinations of $u_1, \ldots, u_s$, the ideal $\mathfrak{b} = (y_1, \ldots, y_d)$ is a reduction of $q$ and
$\{y_1, \ldots, y_d\}$ is a system of parameters of $A$.
\end{theorem}

\begin{proof}
If $d = 0$ then $q^r = (0)$ for some $r > 0$, hence $(0)$ is a reduction of $q$
so that the result holds. We suppose below that $d > 0$.

\textbf{Step 1.} Set $A/m = k$ and consider the polynomial ring $k[X_1, \ldots, X_s]$ (or
$k[X]$ for short). For a homogeneous form $\varphi(X) = \varphi(X_1, \ldots, X_s) \in A[X]$
of degree $n$, we write $\bar{\varphi}(X) \in k[X]$ for the polynomial obtained by
reducing the coefficients of $\varphi$ modulo $m$. As in the proof of Theorem 5 we say that
$\bar{\varphi}(X) \in k[X]$ is a null-form of $q$ if $\varphi(u_1, \ldots, u_s) \in q^{n+1}m$; this notion depends
not just on $q$, but also on $u_1, \ldots, u_s$. However, for fixed $\bar{\varphi}$ it does not
depend on the choice of $\varphi$. We write $Q$ for the ideal of $k[X]$ generated by
all the null-forms of $q$, and call $Q$ the \emph{ideal of null-forms} of $q$. One sees
easily that all the homogeneous elements of $Q$ are null-forms of $q$, and
that the graded ring $k[X]/Q$ has graded component of degree $n$ isomorphic
to $q^n/q^{n+1}m$, so that we have
\[
  k[X]/Q = \bigoplus_{n \ge 0} q^n/q^{n+1}m = \gr_q(A) \otimes_A k.
\]
Write $\varphi(n)$ for the Hilbert function of $k[X]/Q$; then
\[
  l(q^n/q^{n+1}) \le \varphi(n) = l(q^n/q^{n+1}m) \le l(A/q) \cdot l(q^n/q^{n+1})
\]
(see the proof of Theorem 5). We know that for $n \gg 0$, the function $l(q^n/q^{n+1})$
is a polynomial in $n$ of degree $d - 1$ (where $d = \dim A$). Thus from the
above inequality, $\varphi$ is also a polynomial of degree $d - 1$, so that by
Theorem~\ref{thm:13.8}, (ii), we have $\dim k[X]/Q = d$.

Now set $V = \bigoplus kX_i$, and let $P_1, \ldots, P_t$ be the minimal prime divisors of
$Q$. By the assumption that $d > 0$, we have $P_i \nsupseteq V$, so that $P_i \cap V$ is a
proper vector subspace of $V$. Since $k$ is an infinite field,
\[
  V \neq \bigcup_{i=1}^{t} (V \cap P_i).
\]
Hence we can take a linear form $l_1(X) \in V$ not belonging to any $P_i$. If
$d \ge 1$ then similarly we can take $l_2(X) \in V$ such that $l_2(X)$ is not contained
in any minimal prime divisor of $(Q, l_1(X))$, and, proceeding in the same
way, we get $l_1(X), \ldots, l_d(X) \in V$ such that $(Q, l_1, \ldots, l_d)$ is a primary ideal
belonging to $(X_1, \ldots, X_s)$.

\textbf{Step 2.} We let $\mathfrak{b}$ be the ideal of $A$ generated by linear combinations
\[
  L_i(u) = \sum_{j=1}^{s} a_{ij}u_j \quad \text{(for } 1 \le i \le d\text{)}
\]
of $u_1, \ldots, u_s$ with coefficients in $A$. Then if we
put $l_i(X) = \sum a_{ij}X_j$, a necessary and sufficient condition for $\mathfrak{b}$ to be
a reduction of $q$ is that the ideal $(Q, l_1, \ldots, l_d)$ of $k[X]$ is $(X_1, \ldots, X_s)$-
primary.

\emph{Proof of necessity.} Suppose that $\mathfrak{b}q^r = q^{r+1}$. Then if $M(X)$ is a
monomial of degree $r + 1$ in $X_1, \ldots, X_s$, we can write
\[
  M(u) - \sum_i L_i(u)F_i(u) \in q^{r+2}m
\]
where the $F_i(X)$ are homogeneous forms of degrees $r$ with coefficients in
$A$. Thus
\[
  \bar{M}(X) - \sum_i l_i(X)\bar{F}_i(X) \in Q.
\]
Hence
\[
  (X_1, \ldots, X_s)^{r+1} \subset (Q, l_1, \ldots, l_d).
\]

\emph{Proof of sufficiency.} We go through the same argument in reverse: if
$\bar{M}(X) - \sum l_i\bar{F}_i \in Q$ then
\[
  M(u) - \sum L_i(u)F_i(u) \in q^{r+2}m,
\]
so that $q^{r+1} \subset \mathfrak{b}q^r + q^{r+2}m$; thus by NAK, $q^{r+1} = \mathfrak{b}q^r$.

\textbf{Step 3.} Putting together Steps 1 and 2 we see that $q$ has a reduction
$\mathfrak{b} = (y_1, \ldots, y_d)$ generated by $d$ elements. Both $q$ and its reduction $\mathfrak{b}$ are
$m$-primary ideals, so that $y_1, \ldots, y_d$ is a system of parameters of $A$. We
are going to prove that there exists a finite number of polynomials $D_\alpha(Z_{ij})$
for $1 \le \alpha \le \mu$ in $sd$ indeterminates $Z_{ij}$ (for $1 \le i \le d$ and $1 \le j \le s$) such that
$d$ linear combinations $y_i = \sum a_{ij}u_j$ (for $1 \le i \le d$) generate a reduction ideal
of $q$ if and only if at least one of $D_\alpha(a_{ij}) \neq 0$. (The expression `$d$ sufficiently
general linear combinations' in the statement of the theorem is quite vague,
but in the present case it has a precise interpretation as above.)

Let $G_1(X), \ldots, G_t(X)$ be a set of generators of $Q$, with $G_j$ homogeneous
of degree $e_j$. For any $sd$ elements $a_{ij}$ of $k$ (for $1 \le i \le d$ and $1 \le j \le s$), set
$l_i(X) = \sum a_{ij}X_j$. We write $I_n$ for the homogeneous component of degree $n$
of a homogeneous ideal $I \subset k[X_1, \ldots, X_s]$, and in particular we write
$(X_1, \ldots, X_s)_n = V_n$, so that
\[
  (X_1, \ldots, X_s)_n \subset (Q, l_1, \ldots, l_d) \iff V_n = (Q, l_1, \ldots, l_d)_n.
\]
Set $c_n = \dim_k V_n$. We have
\[
  (Q, l_1, \ldots, l_d)_n = \left\{\sum l_iF_i + \sum G_jH_j \mid F_i \in V_{n-1} \text{ and } H_j \in V_{n-e_j}\right\}.
\]
Let $K_1, \ldots, K_w$ be the elements obtained as $\sum l_iF_i + \sum G_jH_j$ as the $F_i$ run
through a basis of $V_{n-1}$ and the $H_j$ run independently through a basis
of $V_{n-e_j}$; it is clear that they span $(Q, l_1, \ldots, l_d)_n$. Each of $K_1, \ldots, K_w$ is a
linear combination of the $c_n$ monomials of degree $n$ in the $X_i$, with linear
functions in the $a_{ij}$ as coefficients; we write out these coefficients in a
$c_n \times w$ matrix. If $\psi_{\nu n}(a_{ij})$ for $1 \le \nu \le \rho$ are the $c_n \times c_n$ minors of this matrix
then the necessary and sufficient condition for $(X_1, \ldots, X_s)_n \subset (Q, l_1, \ldots, l_d)$ to
hold is that at least one of the $\psi_{\nu n}(a_{ij})$ is non-zero. Therefore the ideal
$(Q, l_1, \ldots, l_d)$ will fail to be $(X_1, \ldots, X_s)$-primary if and only if the
quantities $a_{ij}$ satisfy $\psi_{\nu n}(a_{ij}) = 0$ for all $n$ and all $\nu$. However, the ring $k[Z_{ij}]$
is Noetherian, so that the ideal of $k[Z_{ij}]$ generated by all of the $\psi_{\nu n}(Z_{ij})$ is
generated by finitely many elements $D_\alpha(Z_{ij})$ for $1 \le \alpha \le \mu$. These $D_\alpha$ clearly
meet our requirements.
\end{proof}

\begin{remark}
The polynomials $D_\alpha(Z_{ij})$ obtained above are in fact the necessary
and sufficient conditions on the coefficients $a_{ij}$ for the system of homo-
geneous equations $l_1(X) = \cdots = l_d(X) = G_1(X) = \cdots = G_t(X) = 0$ to have a
non-trivial solution, and as such they are known as a \emph{system of resultants}.
Here we have avoided appealing to the classical theory of resultants by
following a method given in Shafarevich [Sh].
\end{remark}

If $k = A/m$ is a finite field then Theorem 14 cannot to be used as it
stands, but we can use the following trick. Let $x$ be an indeterminate over
$A$, and set $S = A[x] - m[x]$; then $S$ consists of polynomials having a unit
of $A$ among their coefficients, and so the composite of the canonical maps
$A \to A[x] \to A[x]_S$ is injective. (In fact $S$ does not contain any zero-
divisors of $A[x]$, so that $A \subset A[x] \subset A[x]_S$; for this see [AM], Chap. 1,
Ex. 2.) Following Nagata~\cite{N1} we write $A(x)$ for $A[x]_S$. This is a
Noetherian local ring containing $A$, with maximal ideal $mA(x)$, and the
residue class field $A(x)/mA(x)$ is the field of fractions of $A[x]/m[x] = k[x]$,
that is, the field $k(x)$ of rational functions over $k$; this is an infinite field.
If $q$ is an $m$-primary ideal of $A$ then $qA(x)$ is a primary ideal belonging
to $mA(x)$. Moreover, since $A(x)$ is flat over $A$, we see that quite generally
if $I \supset I'$ are ideals of $A$ such that $I/I' \cong k$, then
\[
  IA(x)/I'A(x) \cong (I/I') \otimes_A A(x) \cong k \otimes_A A(x) = A(x)/mA(x).
\]
This gives $l_A(A/q^n) = l_{A(x)}(A(x)/q^nA(x))$, so that
\[
  \dim A = \dim A(x) \quad \text{and} \quad e(q) = e(qA(x)).
\]
Thus there are many instances when we can discuss properties of
$e(q)$ in terms of $A(x)$, to which Theorem 14 applies.

% ============ 习题 ============
\section*{Exercises to \S\ref{sec:14}}
\emph{Prove the following propositions.}

\begin{exercise}\label{ex:14.1-ex}
Let $(A, m)$ be a Noetherian local ring and set $G = \gr_m(A)$.
\begin{enumerate}[(i)]
  \item If $G$ is an integral domain then so is $A$ (hence Theorem 3 also follows
        from Theorem 4).
  \item Let $k$ be a field, and $A = k[[X, Y]]/(Y^2 - X^3)$; then $A$ is an integral
        domain, but $G$ has nilpotents.
\end{enumerate}
\end{exercise}

\begin{exercise}\label{ex:14.2-ex}
Let $(A, m)$ and $G$ be as above. For $a \in A$, suppose that $a \in m^i$ but $a \notin m^{i+1}$,
and write $a^*$ for the image of $a$ in $m^i/m^{i+1}$, viewed as an element of $G$;
define $a^*$ to be the \emph{leading term} of $a$. Set $0^* = 0$. Then
\begin{enumerate}[(i)]
  \item if $a^*b^* \neq 0$ then $(ab)^* = a^*b^*$;
  \item if $a^*$ and $b^*$ have the same degree and $a^* + b^* \neq 0$ then $(a + b)^* =
        a^* + b^*$;
  \item let $I \subset m$ be an ideal of $A$. Write $I^*$ for the ideal of $G$ generated by all
        the leading terms of elements of $I$; then setting $B = A/I$ and $n = m/I$, we
        have $\gr_n(B) = G/I^*$.
\end{enumerate}
\end{exercise}

\begin{exercise}\label{ex:14.3-ex}
In the above notation, if $G$ is an integral domain and $I = aA$ then
$I^* = a^*G$. If $I = (a_1, \ldots, a_r)$ with $r > 1$ then it can happen that
$I^* \neq (a_1^*, \ldots, a_r^*)$. Construct an example.
\end{exercise}

\begin{exercise}\label{ex:14.4-ex}
Let $(A, m)$ be a regular local ring, and $K$ its field of fractions.
\begin{enumerate}[(i)]
  \item For $0 \neq a \in A$, set $v(a) = i$ if $a \in m^i$ but $a \notin m^{i+1}$; then $v$ extends to an
        additive valuation of $K$.
  \item Let $R$ be the valuation ring of $v$; then $R$ is a DVR of $K$ dominating $A$.
        Let $x_1, \ldots, x_d$ be a regular system of parameters of $A$, and set
        $B = A[x_2/x_1, \ldots, x_d/x_1]$ and $P = x_1B$; then $P$ is a prime ideal of $B$ and
        $R = B_P$.
\end{enumerate}
\end{exercise}

\begin{exercise}\label{ex:14.5-ex}
In the above notation, if $0 \neq f \in m$ then $v(f)$ is equal to the multiplicity of
$A/(f)$.
\end{exercise}

\begin{exercise}\label{ex:14.6-ex}
(Associativity formula for multiplicities.) Let $A$ be a $d$-dimensional Noether-
ian local ring, $x_1, \ldots, x_d$ a system of parameters of $A$, $q = (x_1, \ldots, x_d)$, and
for $s \le d$ let $\mathfrak{a} = (x_1, \ldots, x_s)$. Write $\Gamma$ for the set of all prime divisors of $\mathfrak{a}$
satisfying $\htop p = s$, $\coht p = d - s$. Let $M$ be a finite $A$-module. Use Lech's
lemma to prove the following formula:
\[
  e(q, M) = \sum_{p \in \Gamma} e(q + p/p, M_p),
\]
(in particular, it follows that $\Gamma \neq \emptyset$).
\end{exercise}

\begin{remark}
The name of the formula comes from its connection with the associativity
of intersection product in algebraic geometry. For details, see [S3], pp. 84--5.
\end{remark}

\begin{exercise}\label{ex:14.7-ex}
Let $(A, m)$ be an $n$-dimensional Noetherian local integral domain, with
$n > 1$. If $0 \neq f \in A$ then $A_f$ is a Jacobson ring (see p. 34).
\end{exercise}

% ============ §15 The dimension of extension rings (原书 pp.116-122) ============
\section{The dimension of extension rings}\label{sec:15}\index{dimension of extension rings, section 15}

\subsection*{1. Fibres}

Let $\varphi\colon A \to B$ be a ring homomorphism, and for $p \in \Spec A$, write $\kappa(p) =
A_p/pA_p$; then $\Spec(B \otimes \kappa(p))$ is called the \emph{fibre}\index{fibre} of $\varphi$ over $p$. As we
saw in \S\ref{sec:7}, it can be identified with the inverse image in $\Spec B$ of $p$ under the
map ${}^a\varphi\colon \Spec B \to \Spec A$ induced by $\varphi$. The ring $B \otimes \kappa(p)$ will be called the
\emph{fibre ring} over $p$. When $(A, m)$ is a local ring, $m$ is the unique closed point of
$\Spec A$, and so the spectrum of $B \otimes \kappa(m) = B/mB$ is called the \emph{closed fibre} of
$\varphi$. If $A$ is an integral domain and $K$ its field of fractions then the spectrum of
$B \otimes_A K = B \otimes_A \kappa(0)$ is called the \emph{generic fibre} of $\varphi$.

\begin{theorem}\label{thm:15.1}
Let $\varphi\colon A \to B$ be a homomorphism of Noetherian rings,
and $P$ a prime ideal of $B$; then setting $p = P \cap A$, we have
\begin{enumerate}[(i)]
  \item $\htop P \le \htop p + \dim B_P/pB_P$;
  \item if $\varphi$ is flat, or more generally if the going-down theorem holds between
        $A$ and $B$, then equality holds in (i).
\end{enumerate}
\end{theorem}

\begin{proof}
We can replace $A$ and $B$ by $A_p$ and $B_P$, and assume that $(A, m)$ and
$(B, n)$ are local rings, with $mB \subset n$. Rewriting (i) in the form
\[
  \dim B \le \dim A + \dim B/mB
\]
makes clear the geometrical content. To prove this, take a system of
parameters $x_1, \ldots, x_r$ of $A$, and choose $y_1, \ldots, y_s \in B$ such that their images
in $B/mB$ form a system of parameters of $B/mB$. Then for $\nu, \rho$ large enough
we have $n^\nu \subset mB + \sum y_iB$ and $m^\rho \subset \sum x_jA$, giving $n^{\nu\rho} \subset \sum y_iB + \sum x_jB$.
Hence $\dim B \le r + s$.

(ii) Let $\dim B/mB = s$, and let $n = P_0 \supset P_1 \supset \cdots \supset P_s$ be a strictly decreas-
ing chain of prime ideals of $B$ between $n$ and $mB$. Obviously we have
$P_i \cap A = m$ for $0 \le i \le s$. Now set $\dim A = r$ and let $m = p_0 \supset p_1 \supset \cdots \supset p_r$
be a strictly decreasing chain of prime ideals of $A$; by the going-down
theorem, we can construct a strictly decreasing chain of prime ideals of $B$
\[
  P_s \supset P_{s+1} \supset \cdots \supset P_{s+r} \quad \text{such that } P_{s+i} \cap A = p_i.
\]
Thus $\dim B \ge r + s$, and putting this together with (i) gives equality.
\end{proof}

\begin{theorem}\label{thm:15.2}
Let $\varphi\colon A \to B$ be a homomorphism of Noetherian rings,
and suppose that the going-up theorem holds between $A$ and $B$. Then if
$p$ and $q$ are prime ideals of $A$ such that $p \supset q$, we have
\[
  \dim B \otimes \kappa(p) \ge \dim B \otimes \kappa(q).
\]
\end{theorem}

\begin{proof}
Set $r = \dim B \otimes \kappa(q)$ and $s = \htop(p/q)$. We choose a strictly
increasing chain $Q_0 \subset Q_1 \subset \cdots \subset Q_r$ of prime ideals of $B$ lying over $q$ and
a strictly increasing chain $q = p_0 \subset p_1 \subset \cdots \subset p_s = p$ of prime ideals of $A$.
By the going-up theorem there exists a chain $Q_r \subset Q_{r+1} \subset \cdots \subset Q_{r+s}$ of
prime ideals of $B$ such that $Q_{r+i} \cap A = p_i$. We set $P = Q_{r+s}$; then
\[
  \htop(P/qB) \ge r + s \quad \text{and} \quad P \cap A = p.
\]
Thus applying the previous theorem to the homomorphism $A/q \to B/qB$
induced by $\varphi$ we get $r + s \le \htop(P/qB) \le s + \dim B_P/pB_P$, and therefore
\[
  r \le \dim B_P/pB_P \le \dim B \otimes \kappa(p). \qquad\blacksquare
\]
\end{proof}

\begin{theorem}\label{thm:15.3}
Let $\varphi\colon A \to B$ be a homomorphism of Noetherian rings, and
suppose that the going-down theorem holds between $A$ and $B$. If $p$ and $q$
are prime ideals of $A$ with $p \supset q$ then
\[
  \dim B \otimes \kappa(p) \le \dim B \otimes \kappa(q).
\]
\end{theorem}

\begin{proof}
We may assume that $\htop(p/q) = 1$, and it is enough to prove that,
given a chain $P_0 \subset P_1 \subset \cdots \subset P_r$ of prime ideals of $B$ lying over $p$ such that
$\htop(P_i/P_{i-1}) = 1$ we can construct a chain of prime ideals $Q_0 \subset Q_1 \subset \cdots \subset Q_r$
of $B$ lying over $q$ such that
\[
  Q_i \subset P_i \quad (0 \le i \le r) \quad \text{and} \quad \htop(Q_i/Q_{i-1}) = 1 \quad (0 < i \le r).
\]
We can find $Q_0$ by going down. If $r \ge 1$ then take $x \in p - q$ and let $T_1, \ldots, T_s$
be the minimal prime divisors of $Q_0 + xB$. Then $\htop(T_j/Q_0) = 1$, while
$\htop(P_r/Q_0) \ge 2$, hence we can choose
\[
  y \in P_r - \bigcup_{j=1}^{s} T_j.
\]
Let $Q_1$ be a minimal prime divisor of $Q_0 + yB$ contained in $P_r$. Then
$\htop(Q_1/Q_0) = 1$, and $y \notin T_i$ for all $i$, hence $x \notin Q_1$.
Therefore $Q_1 \cap A \neq p$, and since $\htop(p/q) = 1$ we must have $Q_1 \cap A = q$. By
the same method we can successively construct $Q_1, Q_2, \ldots, Q_r$.
\end{proof}

\subsection*{2. Polynomial and formal power series rings}

\begin{theorem}\label{thm:15.4}
Let $A$ be a Noetherian ring, and $X_1, \ldots, X_n$ indeterminates
over $A$. Then
\[
  \dim A[X_1, \ldots, X_n] = \dim A[[X_1, \ldots, X_n]] = \dim A + n.
\]
\end{theorem}

\begin{proof}
It is enough to consider the case $n = 1$. For any $p \in \Spec A$, the
fibre $A[X] \otimes_A \kappa(p) = \kappa(p)[X]$ is a principal ideal ring, and therefore one-
dimensional; also $A[X]$ is a free $A$-module, hence faithfully flat, so by
Theorem 1, (ii), $\dim A[X] = \dim A + 1$.

For $A[[X]]$ it is not true in general that $A[[X]] \otimes_A \kappa(p)$ and $\kappa(p)[[X]]$
coincide; however, if $m$ is a maximal ideal of $A$ we have
\[
  A[[X]] \otimes_A \kappa(m) = A[[X]] \otimes_A (A/m) = (A/m)[[X]],
\]
and this fibre ring is one-dimensional. Also, as we saw on p. 4, every
maximal ideal $\mathfrak{M}$ of $A[[X]]$ is of the form $\mathfrak{M} = (m, X)$, where $m = \mathfrak{M} \cap A$
is a maximal ideal of $A$. Thus for a maximal ideal $\mathfrak{M}$ of $A[[X]]$ we have
\[
  \htop \mathfrak{M} = \htop(m \cap A) + 1;
\]
conversely, if $m$ is a maximal ideal of $A$ then $\htop(m, X) = \htop m + 1$, and
putting these together gives $\dim A[[X]] = \dim A + 1$.
\end{proof}

\begin{remark}
It is not necessarily true that a maximal ideal of $A[X]$ lies
over a maximal ideal of $A$. For example, if $A$ is a DVR and $t$ a
uniformising element then $A[t^{-1}] = K$ is the field of fractions of $A$, so
that $A[X]/(tX - 1) \cong K$, and $(tX - 1)$ is a maximal ideal of $A[X]$; however,
$(tX - 1) \cap A = (0)$.
\end{remark}

\begin{remark}
It is quite common for fibre rings of $A \to A[[X_1, \ldots, X_n]]$
to have dimension strictly greater than $n$. For example, let $k$ be a field and
set $A = k[Y, Z]$. It is well-known that the field of fractions of $k[[X]]$ has
infinite transcendence degree over $k(X)$ (see [ZS], vol. II, p. 220). Let $u(X)$,
$v(X) \in k[[X]]$ be two elements algebraically independent over $k(X)$, and
define a $k$-homomorphism (continuous for the $X$-adic topology)
\[
  \varphi\colon A[[X]] = k[Y, Z][[X]] \to k[[X]]
\]
by $\varphi(X) = X$, $\varphi(Y) = u(X)$, $\varphi(Z) = v(X)$. If we set $\Ker \varphi = P$ then $P \cap A = (0)$,
and $A[[X]]/P \cong k[[X]]$ is one-dimensional. Now every maximal ideal of
$A[[X]]$ has height $3$, and, as we will see later, $A[[X]]$ is catenary, so that
$\htop P = 2$. Thus we see that the generic fibre of $A \to A[[X]]$ is two-
dimensional.
\end{remark}

\subsection*{3. The dimension inequality}

We say that a ring $A$ is \emph{universally catenary}\index{universally catenary ring} if $A$ is Noetherian and every
finitely generated $A$-algebra is catenary. Since any $A$-algebra generated
by $n$ elements is a quotient of $A[X_1, \ldots, X_n]$, and since a quotient of a
catenary ring is again catenary, a necessary and sufficient condition for a
Noetherian ring $A$ to be universally catenary is that $A[X_1, \ldots, X_n]$ is
catenary for every $n \ge 0$. (In fact it is known that it is sufficient for $A[X]$
to be catenary, compare Theorem 31.7.)

\begin{theorem}[I. S. Cohen~\cite{Coh3}]\label{thm:15.5}
Let $A$ be a Noetherian integral domain, and
$B$ an extension ring of $A$ which is an integral domain. Let $P \in \Spec B$ and
$p = P \cap A$; then we have
\[
  (*) \quad \htop P + \trdeg_{\kappa(p)} \kappa(P) \le \htop p + \trdeg_K B,
\]
where $\trdeg_K B$ is the transcendence degree of the field of quotients of $B$
over that of $A$.
\end{theorem}

\begin{proof}
We may assume that $B$ is finitely generated over $A$. For if the right
hand side is finite and $m$ and $t$ are non-negative integers such that $m < \htop P$
and $t < \trdeg_{\kappa(p)}\kappa(P)$, then there is a prime ideal chain
$P = P_0 \supset P_1 \supset \cdots \supset P_m$ in $B$. Take $a_i \in P_i - P_{i+1}$, $0 \le i < m$, and let
$c_1, \ldots, c_t \in B$ be such that their images modulo $P$ are algebraically
independent over $A/p$. Set $C = A[\{a_i\}, \{c_j\}]$. If the theorem holds for $C$,
then we have $m + t < \htop p + \trdeg_K C \le \htop p + \trdeg_K B$. Letting $m$ and $t$ vary
we see the validity of $(*)$.

We may furthermore assume, by induction, that $B$ is generated over $A$ by
a single element: $B = A[x]$. We can replace $A$ by $A_p$ and $B$ by $B_P = A_p[x]$,
and hence assume that $A$ is local and $p$ its maximal ideal. Set $k = A/p$ and
write $B = A[X]/Q$. If $Q = (0)$ then $B = A[X]$ and by Theorem 1 we have
$\htop P = \htop p + \htop(P/pB)$, and since $B/pB = k[X]$ we have either $P = pB$ or
$\htop(P/pB) = 1$. In both cases the equality holds in $(*)$.

If $Q \neq (0)$ then $\trdeg_K B = 0$. Since $A$ is a subring of $B$ we have
$Q \cap A = (0)$, so that writing $K$ for the field of fractions of $A$ we have
$\htop Q = \htop QK[X] = 1$. Let $P^*$ be the inverse image of $P$ in $A[X]$. Then
$P = P^*/Q$, $\kappa(P) = \kappa(P^*)$, and
\begin{align*}
  \htop P &\le \htop P^* - \htop Q = \htop P^* - 1 \\
          &= \htop p + 1 - \trdeg_{\kappa(p)}\kappa(P^*) - 1 = \htop p - \trdeg_{\kappa(p)}\kappa(P).
\end{align*}
\end{proof}

\begin{definition}
Suppose that $A$ and $B$ satisfy the conditions of the previous
theorem. We refer to $(*)$ as the \emph{dimension inequality}\index{dimension inequality}, and if the equality in $(*)$
holds for every $P \in \Spec B$, we say that the \emph{dimension formula}\index{dimension formula} holds between
$A$ and $B$. The above proof shows that dimension formula holds between $A$
and $A[x_1, \ldots, x_n]$.
\end{definition}

\begin{theorem}[Ratliff]\label{thm:15.6}
A Noetherian ring $A$ is universally catenary if and
only if the dimension formula holds between $A/p$ and $B$ for every prime
ideal $p$ of $A$ and every finitely generated extension ring $B$ of $A/p$ which
is an integral domain.
\end{theorem}

\begin{proof}[Proof of `only if']
If $A$ is universally catenary then so is $A/p$, so that we
need only consider the case that $A$ is an integral domain, and $B$ is a finitely
generated extension ring which is an integral domain. If
$B = A[X_1, \ldots, X_n]/Q$ and $P = P^*/Q$, then since $A[X_1, \ldots, X_n]$ is catenary
we have $\htop P = \htop P^* - \htop Q$, and an easy calculation proves our assertion.
\end{proof}

\begin{proof}[Proof of `if']
We suppose that $A$ is not universally catenary, so that there
exists a finitely generated $A$-algebra $B$ which is not catenary; without loss
of generality we can assume that $B$ is an integral domain. Write $p$ for the
kernel of the homomorphism $A \to B$. There exist prime ideals $P$ and $Q$
of $B$ such that
\[
  P \subset Q, \quad \htop(Q/P) = d \quad \text{but} \quad \htop Q > \htop P + d.
\]
We write $h = \htop P$, choose $a_1, \ldots, a_h \in P$ such that $\htop(a_1, \ldots, a_h) = h$, and
set $I = (a_1, \ldots, a_h)$, so that $P$ is a minimal prime divisor of $I$. Let
\[
  I = q_1 \cap \cdots \cap q_r
\]
be a shortest primary decomposition of $I$, with $P$ the minimal prime
divisor of $q_1$. Then for $b \in Q \cap q_2 \cap \cdots \cap q_r - P$ we have
\[
  I:b^\nu B = q_1 \quad \text{for } \nu = 1, 2, \ldots.
\]
We set $y_i = a_i/b$ for $1 \le i \le h$,
\[
  C = A[y_1, \ldots, y_h], \quad J = (y_1, \ldots, y_h)C \quad \text{and} \quad M = J + QC = J + Q.
\]
Every element of $C$ can be written in the form $u/b^k$ for suitable $k$, with
$u \in (I + bB)^k$, so that if $z \in J \cap B$ then $zb^\nu \in I$ holds for sufficiently large
$\nu$. Hence $z \in I:b^\nu = q_1$. The converse inclusion $q_1 \subset J \cap B$ is obvious, hence
$J \cap B = q_1$. Thus
\[
  M \cap B = (J + Q) \cap B = (J \cap B) + Q = Q,
\]
\[
  C/J \cong B/q_1, \quad \text{and} \quad C/M \cong B/Q.
\]
Therefore, $C_M/JC_M = B_Q/q_1B_Q$ is a $d$-dimensional local ring, and $J$ is
generated by $h$ elements, so that
\[
  \htop M = \dim C_M \le h + d \le \htop Q.
\]
Now $C$ and $B$ have the same field of fractions, and $\kappa(M) = \kappa(Q)$, so that
this inequality implies that the dimension formula does not hold between
$B$ and $C$. This is a contradiction, since we are assuming that the dimension
formula holds between $A/p$ and $B$ and between $A/p$ and $C$, and one sees
easily that it must then hold between $B$ and $C$.
\end{proof}

\subsection*{4. The Rees ring and gr${}_I$(A)}

Let $A$ be a ring, $I$ an ideal of $A$ and $t$ an indeterminate over $A$. Consider
$A[t]$ as a graded ring in the usual way. We obtain a graded ring $R_+ \subset A[t]$
by setting
\[
  R_+ = R_+(A, I) = \left\{\sum c_nt^n \mid c_n \in I^n\right\} = \bigoplus_{n \ge 0} I^nt^n \subset A[t].
\]
If $I = (a_1, \ldots, a_r)$ then $R_+$ can be written $R_+ = A[a_1t, \ldots, a_rt]$, so that $R_+$ is
Noetherian if $A$ is.

$R_+$ is related to the graded ring $\gr_I(A)$ associated with $A$ and $I$ by the fact
that
\[
  \gr_I(A) = \bigoplus_{n \ge 0} I^n/I^{n+1} \cong R_+/IR_+.
\]

Now let $u = t^{-1}$, and consider $A[t, u] = A[t, t^{-1}]$ as a $\Z$-graded ring in
the obvious way. The \emph{Rees ring}\index{Rees ring} $R(A, I)$ is the graded subring
\[
  R = R(A, I) = R_+[u] = \left\{\sum c_nt^n \in A[t, t^{-1}] \;\middle|\;
    \begin{array}{l}
      c_n \in I^n \text{ for } n < 0 \\
      c_n \in A \text{ for } n > 0
    \end{array} \right\}.
\]
Since
\[
  uR = \left\{\sum c_nt^n \;\middle|\; c_n \in I^{n+1} \text{ for } n \ge 0, \; c_n \in A \text{ for } n < 0\right\},
\]
we have $\gr_I(A) = R/uR$.

\smallskip
\noindent Set $S = \{1, u, u^2, \ldots\}$. Then $R_S = R[u^{-1}] = R[t] = A[t^{-1}, t]$, and
$R_S/(1 - u)R_S = A[t^{-1}, t]/(1 - t) = A$. But $R_S/(1 - u)R_S = (R/(1 - u)R)_{\bar{S}}$,
where $\bar{S}$ is the image of $S$ in $R/(1 - u)R$, and since $\bar{S} = 1$, we see that
$R_{\bar{S}}/(1 - u)R_{\bar{S}} = R/(1 - u)R$. Thus we have
\[
  R/(1 - u)R = A \quad \text{and} \quad R/uR = \gr_I(A),
\]
so that the graded ring $\gr_I(A)$ is a `deformation' of the original ring $A$,
with $R$ as `total space of the deformation', in the sense that $R$ contains a
parameter $u$ such that the values $u = 1$ and $0$ correspond to $A$ and $\gr_I(A)$,
respectively.

We also have
\[
  I^nR \cap A = I^n \quad \text{for all } n \ge 0
\]
and this property is often used to reduce problems about powers of $I$ to
the corresponding problems for powers of the principal ideal $uR$.

We conclude this section by applying the dimension inequality to the
study of the dimension of the Rees ring and $\gr_I(A)$.

Let $A$ be a Noetherian ring, $I = \sum_{i=1}^{r} a_iA$ a proper ideal of $A$, and $t$ an
indeterminate over $A$. We set
\[
  u = t^{-1}, \quad R = R(A, I) = A[u, a_1t, \ldots, a_rt] \quad \text{and} \quad G = \gr_I(A).
\]
We have $R \subset A[t, u]$ and $R/uR \cong G$. For any ideal $\mathfrak{a}$ of $A$, set
\[
  \mathfrak{a}' = \mathfrak{a}A[t, u] \cap R.
\]
Then $\mathfrak{a}' \cap A = \mathfrak{a}A[t, u] \cap A = \mathfrak{a}$, so that for $\mathfrak{a}_1 \neq \mathfrak{a}_2$ we have $\mathfrak{a}_1' \neq \mathfrak{a}_2'$.
Moreover, if $p$ is a prime ideal of $A$ then $p'$ is prime in $R$, and the same
thing goes for primary ideals. If $(0) = q_1 \cap \cdots \cap q_n$ is a primary decompo-
sition of $(0)$ in $A$ then $(0) = q_1' \cap \cdots \cap q_n'$ is a primary decomposition of $(0)$
in $R$. Hence if $p_{0i}$ (for $1 \le i \le m$) are all the minimal prime ideals of $A$ then
$\{p_{0i}'\}_{1 \le i \le m}$ is the set of all minimal prime ideals of $R$. Let $p$ be a prime
ideal of $A$ with $\htop p = h$, and let $p = p_0 \supset p_1 \supset \cdots \supset p_h$ be a strictly decreasing
chain of prime ideals of $A$; then $p' \supset p_1' \supset \cdots \supset p_h'$ is a strictly descending
chain of prime ideals of $R$, so that
\[
  \htop p \le \htop p'.
\]
Conversely, suppose that $P \in \Spec R$ and $P \cap A = p$. Let $p_{0i}$ be a minimal
prime of $R$ contained in $P$ and such that $\htop P = \htop (P/p_{0i}')$; then $R/p_{0i}' \cong
A/p_{0i}$, so that by the dimension inequality
\[
  \htop P = \htop (P/p_{0i}') \le \htop (p/p_{0i}) + 1 - \trdeg_{\kappa(p_{0i})}\kappa(P)
\]
\[
  \le \htop p + 1.
\]
Hence $\dim R \le \dim A + 1$. On the other hand $A[u, t] = R[u^{-1}]$ is a
localisation of $R$ so that $\dim R \ge \dim A[u, t] = \dim A + 1$, so that finally
\[
  \dim R = \dim A + 1.
\]
Moreover, for any $p \in \Spec A$ we set $\bar{a}_i = \text{image of } a_i \bmod p$, so that
\[
  R/p' = (A/p)[u, \bar{a}_1t, \ldots, \bar{a}_rt],
\]
and hence $\trdeg_{\kappa(p)} \kappa(p') = 1$; carrying out the above
calculation using the dimension inequality with $p'$ in place of $P$ we get
$\htop p' \le \htop p$, and so
\[
  \htop p = \htop p'.
\]
We now choose a maximal ideal $m$ of $A$ containing $I$; then since
$R/m' = (A/m)[u]$ we see that $\mathfrak{M} = (m', u)$ is a maximal ideal of $R$ and
$\mathfrak{M} \neq m'$, so that $\htop \mathfrak{M} > \htop m'$. However, by the dimension inequality, we
have $\htop \mathfrak{M} \le \htop m + 1 = \htop m' + 1$. Thus
\[
  \htop \mathfrak{M} = \htop m' + 1 = \htop m + 1.
\]
The element $u$ is a non-zero-divisor of $R$ so that considering a system of
parameters gives $\htop (\mathfrak{M}/uR) = \htop \mathfrak{M} - 1 = \htop m$. Thus providing that there
exists a maximal ideal such that $\htop m = \dim A$ containing $I$, (in particular
if $A$ is local), then we have
\[
  \dim G = \dim (R/uR) = \dim A.
\]
We summarise the above in the following theorem.

\begin{theorem}\label{thm:15.7}
Let $A$ be a Noetherian ring and $I$ a proper ideal; then
setting $R = R(A, I)$ and $G = \gr_I(A)$ we have
\[
  \dim R = \dim A + 1, \quad \dim G \le \dim A.
\]
If in addition $A$ is local, then
\[
  \dim G = \dim A.
\]
\end{theorem}

% ============ 习题 ============
\section*{Exercises to \S\ref{sec:15}}
\emph{Let $k$ be a field.}

\begin{exercise}\label{ex:15.1-ex}
Let $A = k[X, Y] \subset B = k[X, Y, X/Y]$, and $P = (Y, X/Y)B$, $p = (X, Y)A$;
then check that $P \cap A = p$, $\htop P = \htop p = 2$, and $\dim B_P/pB_P = 1$, and hence
that
\[
  \htop P < \htop p + \dim B_P/pB_P.
\]
Show also by a concrete example that the going-down theorem does not
hold between $A$ and $B$.
\end{exercise}

\begin{exercise}\label{ex:15.2-ex}
Does the going-up theorem hold between $A$ and $B$, where $A = k[x] \subset
B = k[x, Y]$?
\end{exercise}

\begin{exercise}\label{ex:15.3-ex}
In Theorem~\ref{thm:15.7}, construct an example where $\dim G < \dim A$.
\end{exercise}

